\documentclass[12pt,a4paper]{article}

%% ============================================================
%% Packages
%% ============================================================
\usepackage[utf8]{inputenc}
\usepackage[T1]{fontenc}
\usepackage{lmodern}
\usepackage{amsmath,amssymb,amsthm}
\usepackage{mathtools}
\usepackage{bm}
\usepackage{enumitem}
\usepackage{booktabs}
\usepackage{array}
\usepackage{hyperref}
\usepackage[margin=1in]{geometry}
\usepackage{cleveref}

%% ============================================================
%% Theorem environments
%% ============================================================
\theoremstyle{plain}
\newtheorem{theorem}{Theorem}[section]
\newtheorem{lemma}[theorem]{Lemma}
\newtheorem{proposition}[theorem]{Proposition}
\newtheorem{corollary}[theorem]{Corollary}

\theoremstyle{definition}
\newtheorem{definition}[theorem]{Definition}
\newtheorem{example}[theorem]{Example}

\theoremstyle{remark}
\newtheorem{remark}[theorem]{Remark}

%% ============================================================
%% Custom commands
%% ============================================================
\newcommand{\Psys}{P_{\mathrm{sys}}}
\newcommand{\Peff}{P_{\mathrm{eff}}}
\newcommand{\Fsys}{F_{\mathrm{sys}}}
\newcommand{\Rel}{\mathrm{Rel}}
\DeclareMathOperator*{\argmax}{arg\,max}

%% ============================================================
%% Title
%% ============================================================
\title{%
  Probabilistic Scaling Laws for Agent Systems:\\[4pt]
  \large From Multiplication Laws to Unified Reliability Polynomials%
}
\author{%
  Agent O \and Agent C \and Orchestrator%
}
\date{March 2026}

\begin{document}
\maketitle

%% ============================================================
%% Abstract
%% ============================================================
\begin{abstract}
We develop a rigorous probabilistic framework for the scaling behavior of
multi-agent systems by modeling each agent as a Bernoulli random variable and
deriving composition laws that govern system-level performance.  Our results
span five regimes: (1)~independent sequential pipelines, where system
reliability decays exponentially with a universal throughput constant
$\Psys(n^*) = e^{-1}$; (2)~dependent pipelines, analyzed as absorbing Markov
chains with coordination gain determined by the geometric mean of coordination
multipliers; (3)~parallel ensembles, exhibiting diminishing returns governed by
the Condorcet jury theorem; (4)~hybrid series-parallel topologies, optimized
via a marginal equalization principle with a unique efficiency maximum; and
(5)~capability saturation, characterized by architecture-specific thresholds
and a phase diagram in capability--cost space.  All five regimes are unified as
special cases of a reliability polynomial on the task as a Directed Acyclic Graph (DAG).  Our theoretical
predictions are quantitatively consistent with the empirical findings of Kim
et~al.\ (2025), including the tool-coordination tradeoff ($R^2 = 0.267$),
capability saturation at $\sim\!45\%$ baseline, centralized coordination
yielding $+81\%$ on structured tasks, independent coordination yielding
$-70\%$ on sequential planning, and task-contingent optimal architecture with
$87\%$ prediction accuracy.
\end{abstract}

\tableofcontents
\newpage

%% ============================================================
%% Section 0: Introduction
%% ============================================================
\section{Introduction}
\label{sec:introduction}

The question of how multi-agent systems scale --- whether adding agents improves
or degrades overall performance --- is fundamental to the design of AI systems
that coordinate multiple language model agents.  Empirical work by Kim
et~al.~(2025) has identified several striking regularities: error amplification
in sequential pipelines, diminishing returns in parallel ensembles, a
capability saturation threshold near $45\%$, and task-contingent optimal
architectures predictable from task topology features.  These findings call
for a formal mathematical framework that can explain and predict scaling
behavior from first principles.

This paper provides such a framework.  We model each agent as a Bernoulli
random variable and derive the probability laws governing system success
under five increasingly general task topologies:

\begin{enumerate}[label=\textbf{Regime \arabic*:}, leftmargin=*]
  \item \textbf{Independent Sequential Pipelines} (\Cref{sec:regime1}).
    The multiplication law yields exponential decay $\Psys = p^n$, a universal
    throughput constant $e^{-1}$, and the result that adding agents always
    degrades effective performance for any coordination cost $c > 0$.

  \item \textbf{Dependent Sequential Pipelines} (\Cref{sec:regime2}).
    The chain rule decomposition, Markov pipeline analysis, and centralized
    validation theory establish when coordination improves reliability and
    when it amplifies errors.

  \item \textbf{Parallel Ensembles} (\Cref{sec:regime3}).
    Diminishing returns follow a geometric sequence with ratio $(1-p)$, the
    Condorcet jury theorem governs majority voting, and parallel architectures
    strictly dominate sequential ones for $n \geq 2$.

  \item \textbf{Hybrid Series-Parallel Topologies} (\Cref{sec:regime4}).
    Recursive composition on series-parallel DAGs, the marginal equalization
    principle for optimal resource allocation, and a unique efficiency maximum
    per topology class.

  \item \textbf{Capability Saturation and the Unified Scaling Law}
    (\Cref{sec:regime5}).
    Architecture-specific saturation thresholds, a phase diagram in $(p,c)$-space,
    and the reliability polynomial as the unifying mathematical object
    encompassing all regimes.
\end{enumerate}

\subsection{Notation}

Throughout, $A_i$ denotes an agent modeled as a Bernoulli random variable with
success probability $p_i = P(A_i = 1)$.  We write $p$ for the uniform
capability when all agents share the same success probability, $q = 1-p$ for
the failure probability, and $c$ for the per-link (or per-agent) coordination
cost.  The system success probability is $\Psys$ and the effective system
probability (including cost) is $\Peff$.  The error amplification factor is
$\alpha = p(1-c)$.

\subsection{Relationship to Prior Work}

Our framework is motivated by the empirical findings of Kim et~al.\ (2025),
\emph{Towards a Science of Scaling Agent Systems}, which studied multi-agent
coordination across diverse benchmarks.  We provide theoretical foundations for
their observations, deriving each empirical pattern from the probabilistic
structure of the task topology.  The correspondence between our predictions and
their findings is summarized in \Cref{sec:regime5} (Theorem~\ref{thm:empirical-correspondence}).

\newpage

%% ============================================================
%% Section 1: Definitions
%% ============================================================
\section{Primitive Definitions}
\label{sec:definitions}

\begin{definition}[Agent Random Variable]
\label{def:agent}
An \emph{agent} is a Bernoulli random variable $A_i : \Omega \to \{0, 1\}$ defined on a
probability space $(\Omega, \mathcal{F}, P)$, where $A_i = 1$ denotes successful completion
of the agent's assigned subtask and $A_i = 0$ denotes failure. The \emph{capability} of
agent $A_i$ is its success probability:
\[
  p_i \;:=\; P(A_i = 1) \;\in\; [0, 1].
\]
When all agents share a common capability we write $p_i = p$ for all $i$ and call $p$
the \emph{uniform capability}.
\end{definition}

\begin{definition}[Sequential Pipeline]
\label{def:pipeline}
A \emph{sequential pipeline} of $n$ agents is an ordered tuple $(A_1, A_2, \ldots, A_n)$
in which the system produces a correct output if and only if every agent succeeds.
The \emph{system indicator} is
\[
  S_n \;:=\; \prod_{i=1}^{n} A_i,
\]
so $S_n = 1$ iff $A_i = 1$ for all $i \in \{1, \ldots, n\}$.
\end{definition}

\begin{definition}[System Success Probability]
\label{def:psys}
The \emph{system success probability} for a pipeline of $n$ agents is
\[
  P_{\mathrm{sys}}(n) \;:=\; P(S_n = 1) \;=\; P\!\Bigl(\,\bigcap_{i=1}^{n} \{A_i = 1\}\Bigr).
\]
When the agents are mutually independent this reduces to
$P_{\mathrm{sys}}(n) = \prod_{i=1}^{n} p_i$.
In the dependent case the chain rule gives
$P_{\mathrm{sys}}(n) = \prod_{i=1}^{n} P(A_i = 1 \mid A_1 = 1, \ldots, A_{i-1} = 1)$.
\end{definition}

\begin{definition}[Coordination Cost Function --- Multiplicative Model]
\label{def:coord-cost}
Let $c \in [0,1)$ be the \emph{per-link coordination cost}. In a sequential pipeline of
$n$ agents there are $(n-1)$ inter-agent links. Each link imposes a multiplicative
overhead factor $(1-c)$ on the system. The \emph{effective system success probability}
under coordination cost is
\[
  P_{\mathrm{eff}}(n) \;:=\; P_{\mathrm{sys}}(n) \cdot (1-c)^{n-1}.
\]
We adopt the multiplicative model (rather than an additive model $p_i \mapsto p_i - c$)
for three reasons:
\begin{enumerate}
  \item It preserves the constraint $P_{\mathrm{eff}}(n) \in [0,1]$ automatically.
  \item It composes naturally with the multiplicative structure of independent success
        probabilities.
  \item The additive model $p - c$ can become negative for $c > p$, requiring artificial
        truncation.
\end{enumerate}
\end{definition}

\begin{definition}[Net Benefit Function]
\label{def:net-benefit}
The \emph{net benefit} of adding the $(n+1)$-th agent to a pipeline is
\[
  B(n) \;:=\; P_{\mathrm{eff}}(n+1) - P_{\mathrm{eff}}(n).
\]
An agent addition is \emph{beneficial} if $B(n) > 0$ and \emph{detrimental} if $B(n) < 0$.
\end{definition}

\begin{definition}[Task Topology]
\label{def:topology}
A \emph{task topology} is a directed acyclic graph $G = (V, E)$ where each vertex
$v \in V$ is an agent and each directed edge $(u, v) \in E$ indicates that agent $v$
receives output from agent $u$. The three primitive topologies are:
\begin{enumerate}
  \item \textbf{Sequential (chain)}: $G$ is a directed path $A_1 \to A_2 \to \cdots \to A_n$.
  \item \textbf{Parallel (ensemble)}: $G$ consists of $n$ agents with a common source and a
        common sink (aggregator), with no edges between agents.
  \item \textbf{Hybrid (DAG)}: $G$ is a general DAG combining sequential and parallel
        substructures.
\end{enumerate}
\end{definition}

\begin{definition}[Error Amplification Factor]
\label{def:error-amp}
For a sequential pipeline of $n$ agents with uniform capability $p$ and per-link
coordination cost $c$, the \emph{system failure probability} is
\[
  F_{\mathrm{sys}}(n) \;=\; 1 - P_{\mathrm{eff}}(n).
\]
The \emph{error amplification factor} $\alpha$ is defined by the relation
\[
  P_{\mathrm{eff}}(n) \;=\; \alpha^{\,n-1} \cdot p,
\]
where $\alpha := p(1-c) \in (0,1)$ for $p < 1$ and $c > 0$. Equivalently, $-\ln \alpha$
is the exponential decay rate of system reliability per additional agent.
\end{definition}

\section{Regime 2 --- Dependent Pipeline Definitions}

\begin{definition}[Dependent Pipeline (Conditional Chain)]
\label{def:dependent-pipeline}
A \emph{dependent pipeline} of $n$ agents is a sequential pipeline $(A_1, \ldots, A_n)$
in which the success of agent $A_i$ may depend on the outcomes of all preceding agents.
The system succeeds if and only if all agents succeed. By the chain rule of probability:
\[
  P_{\mathrm{sys}}(n) \;=\; \prod_{i=1}^{n} P(A_i = 1 \mid A_1 = 1, \ldots, A_{i-1} = 1),
\]
with the convention that $P(A_1 = 1 \mid \emptyset) = P(A_1 = 1) = p_1$.

We write $q_i := P(A_i = 1 \mid A_1 = 1, \ldots, A_{i-1} = 1)$ for the \emph{conditional
success probability} of agent $i$ given that all predecessors succeeded. Then
$P_{\mathrm{sys}}(n) = \prod_{i=1}^{n} q_i$.
\end{definition}

\begin{definition}[Coordination Gain Function]
\label{def:coord-gain}
For a pipeline of $n$ agents with marginal success probabilities $p_1, \ldots, p_n$,
the \emph{coordination gain} is the difference in system success probability between
a dependent (coordinated) pipeline and an independent pipeline:
\[
  \Delta(n) \;:=\; P_{\mathrm{sys}}^{\mathrm{dep}}(n) - P_{\mathrm{sys}}^{\mathrm{indep}}(n)
  \;=\; \prod_{i=1}^{n} q_i \;-\; \prod_{i=1}^{n} p_i.
\]
Coordination is \emph{beneficial} if $\Delta(n) > 0$ and \emph{detrimental} if $\Delta(n) < 0$.
\end{definition}

\begin{definition}[Markov Pipeline]
\label{def:markov-pipeline}
A dependent pipeline is \emph{Markovian} if the conditional success probability of each
agent depends only on the immediately preceding agent's outcome, not on the full history:
\[
  P(A_{i+1} = 1 \mid A_1 = 1, \ldots, A_i = 1) \;=\; P(A_{i+1} = 1 \mid A_i = 1) \;=:\; r_i
\]
for $i = 1, \ldots, n-1$, where $r_i \in [0,1]$ is the \emph{transition probability} from
agent $i$ to agent $i+1$. In a Markov pipeline:
\[
  P_{\mathrm{sys}}(n) \;=\; p_1 \cdot \prod_{i=1}^{n-1} r_i.
\]
When all transition probabilities are equal, $r_i = r$, we write $P_{\mathrm{sys}}(n) = p_1 \cdot r^{n-1}$.
\end{definition}

\begin{definition}[Centralized Coordinator (Hub-and-Spoke)]
\label{def:hub-spoke}
A \emph{centralized coordinator} is a distinguished agent $V$ (the validator) that
mediates the pipeline: each worker agent $A_i$ submits its output to $V$, which
accepts or rejects before passing to $A_{i+1}$. The validator has detection probability
$d \in [0,1]$ (probability of catching an error given one occurred) and false positive
rate $f \in [0,1]$ (probability of rejecting correct output). The effective conditional
success probability under centralized validation is:
\[
  q_i^{\mathrm{val}} \;=\; P(A_i \text{ accepted and correct} \mid \text{predecessors succeeded})
  \;=\; p_i(1 - f) + (1 - p_i)(1 - d) \cdot 0 \;=\; p_i(1 - f).
\]
More precisely, the system succeeds at step $i$ if $A_i$ succeeds \emph{and} the validator
does not falsely reject, giving $q_i^{\mathrm{val}} = p_i(1-f)$. Errors are caught
with probability $d$, preventing downstream propagation.
\end{definition}

% End of definitions


\newpage

%% ============================================================
%% Section 2: Regime 1
%% ============================================================
\section{Regime 1 --- Independent Agents (Multiplication Law)}
\label{sec:regime1}

Throughout this section, let $(A_1, A_2, \ldots, A_n)$ be a sequential pipeline of
agents as in Definition~\ref{def:pipeline}, with $A_i \sim \mathrm{Bernoulli}(p_i)$.
Unless stated otherwise, the agents are assumed \textbf{mutually independent}.

%% ============================================================
\subsection{The Multiplication Law}
%% ============================================================

\begin{theorem}[Multiplication Law for Independent Pipelines]
\label{thm:mult-law}
If agents $A_1, \ldots, A_n$ are mutually independent with success probabilities
$p_1, \ldots, p_n$, then
\[
  P_{\mathrm{sys}}(n) \;=\; \prod_{i=1}^{n} p_i.
\]
\end{theorem}

\begin{proof}
By Definition~\ref{def:psys},
\[
  P_{\mathrm{sys}}(n) = P\!\Bigl(\,\bigcap_{i=1}^{n} \{A_i = 1\}\Bigr).
\]
Since the $A_i$ are mutually independent, the probability of the intersection of
independent events equals the product of their probabilities (this is the defining
property of mutual independence applied to the events $\{A_i = 1\}$):
\[
  P\!\Bigl(\,\bigcap_{i=1}^{n} \{A_i = 1\}\Bigr)
  \;=\; \prod_{i=1}^{n} P(A_i = 1)
  \;=\; \prod_{i=1}^{n} p_i.
  \qedhere
\]
\end{proof}

%% ============================================================
\subsection{Exponential Decay of Reliability}
%% ============================================================

\begin{theorem}[Exponential Decay]
\label{thm:exp-decay}
Let all agents have uniform capability $p \in (0,1)$. Then
\[
  P_{\mathrm{sys}}(n) \;=\; p^n \;=\; e^{\,n \ln p},
\]
which decays exponentially in $n$ at rate $|\ln p| = -\ln p > 0$.

More precisely, for any target reliability $P_{\mathrm{sys}}(n) \geq \delta$ with
$\delta \in (0,1)$, the maximum number of agents is
\[
  n_{\max}(\delta, p) \;=\; \Bigl\lfloor \frac{\ln \delta}{\ln p} \Bigr\rfloor.
\]
\end{theorem}

\begin{proof}
By Theorem~\ref{thm:mult-law} with $p_i = p$ for all $i$,
\[
  P_{\mathrm{sys}}(n) = p^n.
\]
Since $0 < p < 1$, we have $\ln p < 0$, so $p^n = e^{n \ln p}$ is a strictly
decreasing exponential function of $n$. The decay rate per agent is
\[
  -\frac{d}{dn}\ln P_{\mathrm{sys}}(n) = -\ln p = |\ln p| > 0.
\]

For the bound, require $p^n \geq \delta$. Taking logarithms (note $\ln p < 0$, so
the inequality reverses):
\[
  n \ln p \geq \ln \delta
  \;\;\Longleftrightarrow\;\;
  n \leq \frac{\ln \delta}{\ln p}.
\]
Since $n$ must be a positive integer, $n_{\max} = \lfloor \ln\delta / \ln p \rfloor$.
\end{proof}

\begin{corollary}[Half-Life of a Pipeline]
\label{cor:half-life}
The number of agents at which system reliability drops to half the single-agent
capability is
\[
  n_{1/2} \;=\; \Bigl\lfloor \frac{\ln(1/2)}{\ln p} \Bigr\rfloor + 1
  \;=\; \Bigl\lfloor \frac{\ln 2}{|\ln p|} \Bigr\rfloor + 1.
\]
For high-capability agents ($p$ close to 1), using $\ln p \approx -(1-p)$:
\[
  n_{1/2} \;\approx\; \frac{\ln 2}{1-p}.
\]
\end{corollary}

\begin{proof}
We seek the smallest $n$ such that $p^n \leq p/2$ (half the single-agent reliability).
Dividing by $p$, this is $p^{n-1} \leq 1/2$, i.e., $(n-1) \geq \ln(1/2)/\ln p =
\ln 2 / |\ln p|$. The approximation follows from the first-order Taylor expansion
$\ln p = \ln(1-(1-p)) \approx -(1-p)$ for $p$ near 1.
\end{proof}

%% ============================================================
\subsection{Heterogeneous Agents: AM--GM Bound}
%% ============================================================

\begin{lemma}[Uniform Capability Maximizes System Reliability]
\label{lem:am-gm}
Among all pipelines of $n$ independent agents with a fixed sum of capabilities
$\sum_{i=1}^{n} p_i = S$ (where $S \leq n$), the system success probability
$P_{\mathrm{sys}} = \prod p_i$ is maximized when $p_i = S/n$ for all $i$.

Equivalently, for any fixed arithmetic mean $\bar{p} = S/n$:
\[
  \prod_{i=1}^{n} p_i \;\leq\; \bar{p}^{\,n},
\]
with equality iff $p_1 = p_2 = \cdots = p_n$.
\end{lemma}

\begin{proof}
This is the AM--GM inequality applied in logarithmic form. Since $\ln$ is concave,
Jensen's inequality gives
\[
  \frac{1}{n}\sum_{i=1}^{n} \ln p_i \;\leq\; \ln\!\Bigl(\frac{1}{n}\sum_{i=1}^{n} p_i\Bigr) = \ln \bar{p},
\]
hence $\sum \ln p_i \leq n \ln \bar{p}$, i.e., $\prod p_i \leq \bar{p}^n$.
Equality in Jensen's inequality holds iff all $p_i$ are equal.
\end{proof}

%% ============================================================
\subsection{Critical Threshold}
%% ============================================================

\begin{theorem}[Critical Threshold for Beneficial Agent Addition]
\label{thm:critical-threshold}
Consider a uniform-capability pipeline with per-link coordination cost $c \in [0,1)$.
The effective system success probability is (Definition~\ref{def:coord-cost}):
\[
  P_{\mathrm{eff}}(n) \;=\; p^n \cdot (1-c)^{n-1}.
\]

Adding an $(n+1)$-th agent is beneficial (i.e., $P_{\mathrm{eff}}(n+1) > P_{\mathrm{eff}}(n)$) if
and only if
\[
  p(1-c) > 1.
\]

Therefore the \emph{critical threshold} is
\[
  \boxed{p^* \;=\; \frac{1}{1-c}.}
\]

\begin{itemize}
  \item For $p > p^*$: every additional agent \emph{increases} effective system reliability.
  \item For $p < p^*$: every additional agent \emph{decreases} effective system reliability.
  \item For $p = p^*$: $P_{\mathrm{eff}}$ is constant in $n$.
\end{itemize}

Note: Since $p \leq 1$ and $p^* = 1/(1-c)$, the condition $p > p^*$ can only be
satisfied when $c = 0$ and $p = 1$. For any $c > 0$, we have $p^* > 1$, which means
$p < p^*$ always holds --- adding agents to a pipeline with coordination cost always
degrades performance. This is the fundamental bottleneck of sequential architectures.
\end{theorem}

\begin{proof}
Compute the ratio:
\[
  \frac{P_{\mathrm{eff}}(n+1)}{P_{\mathrm{eff}}(n)}
  \;=\; \frac{p^{n+1}(1-c)^{n}}{p^{n}(1-c)^{n-1}}
  \;=\; p(1-c).
\]
This ratio is independent of $n$. We have:
\begin{align*}
  P_{\mathrm{eff}}(n+1) > P_{\mathrm{eff}}(n)
  &\;\Longleftrightarrow\; p(1-c) > 1
  \;\Longleftrightarrow\; p > \frac{1}{1-c} = p^*.
\end{align*}

Since $c \in [0,1)$, we have $1-c \in (0,1]$, hence $p^* = 1/(1-c) \geq 1$.
Because $p \leq 1$ by definition, the condition $p > p^*$ requires $p^* < 1$,
which is impossible. When $c = 0$, $p^* = 1$ and the condition becomes $p > 1$,
never satisfied. Thus for $c > 0$, $p^* > 1 > p$ strictly, and adding agents
always reduces $P_{\mathrm{eff}}$.
\end{proof}

\begin{remark}[Interpretation for Zero Coordination Cost]
\label{rmk:zero-cost}
When $c = 0$ (no coordination overhead), the effective probability reduces to
$P_{\mathrm{eff}}(n) = p^n$, which is still decreasing for $p < 1$. The critical threshold
$p^* = 1$ means that only a perfect agent ($p = 1$) can be added without degradation.
This confirms that the sequential pipeline is inherently fragile: it is the
multiplicative structure itself, not just coordination cost, that drives exponential
decay.
\end{remark}

%% ============================================================
\subsection{Optimal Number of Agents}
%% ============================================================

The preceding theorem shows that in a pure sequential pipeline, adding agents always
hurts when $c > 0$. This means the trivially optimal pipeline length is $n^* = 1$.
However, this result assumes agents perform identical tasks. A more realistic model
assumes that the task \emph{requires} decomposition: a task of complexity $T$ must be
divided among agents, and more agents means each agent handles a simpler subtask,
increasing its per-subtask success probability.

\begin{definition}[Task Decomposition Model]
\label{def:task-decomp}
A task of complexity $T > 0$ is divided among $n$ agents. Each agent handles a subtask
of complexity $T/n$. The per-agent success probability is a decreasing function of
subtask complexity:
\[
  p(n) \;=\; g\!\left(\frac{T}{n}\right),
\]
where $g : [0, \infty) \to [0,1]$ is a decreasing function with $g(0) = 1$ and
$\lim_{x \to \infty} g(x) = 0$. A natural choice is the exponential model:
\[
  p(n) \;=\; e^{-T/n}.
\]
\end{definition}

\begin{theorem}[Optimal Pipeline Length --- Exponential Decomposition]
\label{thm:optimal-n}
Under the exponential task decomposition model $p(n) = e^{-T/n}$ with per-link
coordination cost $c \in (0,1)$, the effective system probability is
\[
  P_{\mathrm{eff}}(n) \;=\; e^{-T} \cdot (1-c)^{n-1}.
\]
This is maximized at $n^* = 1$ for all $T > 0$ and $c > 0$.
\end{theorem}

\begin{proof}
Substitute $p(n) = e^{-T/n}$ into $P_{\mathrm{eff}}$:
\[
  P_{\mathrm{eff}}(n)
  = \bigl(e^{-T/n}\bigr)^n \cdot (1-c)^{n-1}
  = e^{-T} \cdot (1-c)^{n-1}.
\]
The factor $e^{-T}$ is independent of $n$. Since $(1-c) < 1$, the factor $(1-c)^{n-1}$
is strictly decreasing in $n$. Hence $P_{\mathrm{eff}}$ is maximized at $n = 1$, where
$P_{\mathrm{eff}}(1) = e^{-T}$.
\end{proof}

\begin{remark}
Theorem~\ref{thm:optimal-n} reveals a degenerate case: the exponential decomposition
$p(n) = e^{-T/n}$ satisfies $[p(n)]^n = e^{-T}$ identically, so task decomposition
yields \emph{no net benefit} to the multiplicative product. All benefit from simpler
subtasks is exactly canceled by having more agents. Coordination cost then strictly
favors fewer agents. To obtain a non-trivial optimum we need decomposition that
yields a net benefit exceeding coordination cost.
\end{remark}

\begin{definition}[Power-Law Decomposition Model]
\label{def:power-decomp}
A more favorable decomposition model assumes per-agent success probability
\[
  p(n) \;=\; 1 - \left(\frac{T}{n}\right)^{\!\beta}
\]
for some $\beta > 0$ and $T/n < 1$ (i.e., $n > T$). The parameter $\beta$ controls
how rapidly success probability improves as subtask complexity decreases.
\end{definition}

\begin{theorem}[Optimal Pipeline Length --- General Decomposition]
\label{thm:optimal-n-general}
Let $p(n)$ be a decomposition model such that the log-effective probability
\[
  \ell(n) \;:=\; \ln P_{\mathrm{eff}}(n) \;=\; n \ln p(n) + (n-1)\ln(1-c)
\]
is a concave function of $n$ (treating $n$ as continuous). Then the optimal pipeline
length $n^*$ satisfies
\[
  \boxed{\ln p(n^*) + n^* \cdot \frac{p'(n^*)}{p(n^*)} + \ln(1-c) = 0.}
\]
Equivalently,
\[
  n^* \cdot \frac{p'(n^*)}{p(n^*)} \;=\; -\ln p(n^*) - \ln(1-c)
  \;=\; -\ln\!\bigl[p(n^*)(1-c)\bigr].
\]
\end{theorem}

\begin{proof}
Differentiate $\ell(n)$ with respect to $n$ (treating $n$ as continuous):
\[
  \ell'(n) = \ln p(n) + n \cdot \frac{p'(n)}{p(n)} + \ln(1-c).
\]
Setting $\ell'(n^*) = 0$ gives the stated equation. If $\ell$ is concave, this
critical point is a global maximum.
\end{proof}

\begin{corollary}[Explicit Solution for Near-Perfect Agents]
\label{cor:optimal-approx}
When $p$ is close to 1 (high capability), write $p(n) = 1 - \epsilon(n)$ where
$\epsilon(n) \ll 1$. Then $\ln p(n) \approx -\epsilon(n)$ and $p'(n)/p(n) \approx
-\epsilon'(n)$. The optimality condition becomes
\[
  -\epsilon(n^*) - n^* \epsilon'(n^*) + \ln(1-c) = 0.
\]
For the power-law model $\epsilon(n) = (T/n)^\beta$ with $\epsilon'(n) =
-\beta T^\beta / n^{\beta+1}$:
\[
  -\frac{T^\beta}{(n^*)^\beta} + \frac{\beta T^\beta}{(n^*)^\beta} + \ln(1-c) = 0,
\]
giving
\[
  \boxed{n^* \;=\; T \cdot \left(\frac{\beta - 1}{|\ln(1-c)|}\right)^{\!1/\beta}}
  \quad (\text{for } \beta > 1).
\]
For $\beta \leq 1$, no finite optimum exists (the numerator is non-positive), meaning
decomposition benefit is too weak to overcome coordination cost.
\end{corollary}

\begin{proof}
Substituting $\epsilon(n) = (T/n)^\beta$:
\[
  -\frac{T^\beta}{n^\beta} + \beta \cdot \frac{T^\beta}{n^\beta} + \ln(1-c) = 0
  \;\;\Longrightarrow\;\;
  (\beta - 1) \frac{T^\beta}{n^\beta} = -\ln(1-c) = |\ln(1-c)|.
\]
Solving for $n$:
\[
  n^\beta = \frac{(\beta-1) T^\beta}{|\ln(1-c)|}
  \;\;\Longrightarrow\;\;
  n = T \left(\frac{\beta-1}{|\ln(1-c)|}\right)^{1/\beta}.
\]
This requires $\beta > 1$ for the expression to be positive and real.
\end{proof}

%% ============================================================
\subsection{Error Amplification in Independent Pipelines}
%% ============================================================

\begin{theorem}[Error Amplification Factor]
\label{thm:error-amp}
For a uniform-capability pipeline with per-agent success probability $p$ and
coordination cost $c$, define the \emph{amplification factor} $\alpha := p(1-c)$.
Then:
\[
  P_{\mathrm{eff}}(n) = \frac{p}{1-c} \cdot \alpha^n = \frac{p \cdot \alpha^n}{1-c}.
\]
Equivalently, $P_{\mathrm{eff}}(n) = p \cdot \alpha^{n-1}$.

The system failure probability satisfies:
\[
  F_{\mathrm{sys}}(n) = 1 - p \cdot \alpha^{n-1}.
\]

For $\alpha < 1$ (which holds whenever $p < 1$ or $c > 0$ with $p \leq 1$):
\begin{enumerate}
  \item $P_{\mathrm{eff}}(n) \to 0$ exponentially as $n \to \infty$.
  \item $F_{\mathrm{sys}}(n) \to 1$ exponentially as $n \to \infty$.
  \item The \emph{reliability half-life} is $n_{1/2} = 1 + \ln 2 / |\ln \alpha|$.
\end{enumerate}
\end{theorem}

\begin{proof}
From Definition~\ref{def:coord-cost}:
\[
  P_{\mathrm{eff}}(n) = p^n (1-c)^{n-1} = p \cdot [p(1-c)]^{n-1} = p \cdot \alpha^{n-1}.
\]
Since $\alpha = p(1-c)$ and both $p \leq 1$, $(1-c) \leq 1$ with at least one strict,
we have $\alpha < 1$. Therefore $\alpha^{n-1} \to 0$ as $n \to \infty$, giving
conclusions (1) and (2).

For (3), set $P_{\mathrm{eff}}(n_{1/2}) = P_{\mathrm{eff}}(1)/2 = p/2$:
\[
  p \cdot \alpha^{n_{1/2}-1} = p/2
  \;\Longrightarrow\;
  \alpha^{n_{1/2}-1} = 1/2
  \;\Longrightarrow\;
  n_{1/2} = 1 + \frac{\ln 2}{|\ln \alpha|}.
  \qedhere
\]
\end{proof}

%% ============================================================
\subsection{Throughput-Optimal Pipeline Length}
%% ============================================================

\begin{theorem}[Universal Throughput Optimum]
\label{thm:throughput-optimum}
Define the \emph{throughput} of a uniform-capability pipeline as $\Theta(n) := n \cdot P_{\mathrm{sys}}(n) = n \, p^n$,
representing the expected number of successfully completed subtasks. Then:
\begin{enumerate}
  \item The throughput-maximizing pipeline length (treating $n$ as continuous) is
  \[
    n^* \;=\; -\frac{1}{\ln p}.
  \]
  \item At this optimum, the system success probability attains the \emph{universal constant}
  \[
    P_{\mathrm{sys}}(n^*) \;=\; e^{-1} \;\approx\; 0.3679,
  \]
  independent of $p$.
  \item The maximum throughput is $\Theta^* = n^* \cdot e^{-1} = -1/(e \ln p)$.
\end{enumerate}
\end{theorem}

\begin{proof}
For (1), differentiate $\Theta(n) = n p^n$ with respect to $n$:
\[
  \Theta'(n) = p^n + n p^n \ln p = p^n(1 + n \ln p).
\]
Setting $\Theta'(n^*) = 0$ and noting $p^{n^*} > 0$, we get $1 + n^* \ln p = 0$,
hence $n^* = -1/\ln p$. The second derivative $\Theta''(n^*) = p^{n^*} \ln p < 0$
confirms this is a maximum.

For (2), substitute $n^* = -1/\ln p$ into $P_{\mathrm{sys}} = p^n$:
\[
  P_{\mathrm{sys}}(n^*) = p^{-1/\ln p} = e^{\,\ln p \cdot (-1/\ln p)} = e^{-1}.
\]

Statement (3) follows immediately: $\Theta^* = n^* \cdot e^{-1} = -1/(e \ln p)$.
\end{proof}

\begin{remark}
The universality of $P_{\mathrm{sys}}(n^*) = e^{-1}$ means that \emph{every} throughput-optimal
pipeline, regardless of agent capability $p$, operates at the same system success probability.
Higher-capability agents simply allow longer pipelines (larger $n^*$) while maintaining
this fixed reliability. This provides a natural bridge to Regime~2: if dependent agents
can keep $P_{\mathrm{sys}} > e^{-1}$ at the same pipeline length, coordination is provably beneficial.
\end{remark}

%% ============================================================
\subsection{Correspondence with Empirical Findings}
%% ============================================================

\begin{remark}[Empirical Correspondence]
\label{rmk:empirical}
The results of this section correspond to the following empirical findings from
Kim et al.\ (2025):

\begin{enumerate}
\item \textbf{Independent coordination: $-70\%$ on sequential planning tasks.}
  Theorem~\ref{thm:critical-threshold} shows that for any $c > 0$, adding
  independent agents to a sequential pipeline always degrades performance.
  With uniform $p \approx 0.7$ and modest $c \approx 0.05$, a 3-agent pipeline gives
  $P_{\mathrm{eff}}(3) = 0.7^3 \cdot 0.95^2 \approx 0.310$, a 56\% reduction from the
  single-agent $P_{\mathrm{eff}}(1) = 0.7$.

\item \textbf{Error amplification: independent agents amplify errors $\sim\alpha^n$.}
  Theorem~\ref{thm:error-amp} derives this exactly with $\alpha = p(1-c) < 1$.

\item \textbf{Tool-coordination tradeoff: overhead compounds with environment complexity.}
  In the task decomposition model (Theorem~\ref{thm:optimal-n-general}), increasing
  task complexity $T$ shifts the optimal $n^*$ and changes the balance between
  decomposition benefit and coordination cost. The coordination overhead
  $(1-c)^{n-1}$ compounds geometrically with pipeline length.

\item \textbf{Capability saturation at $\sim 45\%$ baseline.}
  For $p = 0.45$ and $c = 0$, a 2-agent pipeline gives $0.45^2 = 0.2025$, already
  below the single-agent baseline. This prefigures the saturation result of Regime~5:
  the critical threshold $p^* = 1/(1-c)$ means sequential coordination is generically
  unprofitable, and the saturation threshold emerges from the balance with parallel
  architectures (Regime~3).
\end{enumerate}
\end{remark}

%% ============================================================
\subsection{Summary of Regime 1 Results}
%% ============================================================

We collect the principal results:

\begin{enumerate}
  \item \textbf{Multiplication Law} (Thm~\ref{thm:mult-law}):
    $P_{\mathrm{sys}} = \prod p_i$.

  \item \textbf{Exponential Decay} (Thm~\ref{thm:exp-decay}):
    $P_{\mathrm{sys}} = p^n = e^{n\ln p}$, decaying at rate $|\ln p|$ per agent.

  \item \textbf{AM--GM Bound} (Lem~\ref{lem:am-gm}):
    Heterogeneous agents are strictly worse than uniform agents with the same mean.

  \item \textbf{Critical Threshold} (Thm~\ref{thm:critical-threshold}):
    $p^* = 1/(1-c)$; for $c > 0$, $p^* > 1$, so adding agents always hurts.

  \item \textbf{Optimal Pipeline Length} (Thm~\ref{thm:optimal-n-general},
    Cor~\ref{cor:optimal-approx}):
    Under power-law decomposition with $\beta > 1$:
    $n^* = T[(\beta-1)/|\ln(1-c)|]^{1/\beta}$.

  \item \textbf{Error Amplification} (Thm~\ref{thm:error-amp}):
    $\alpha = p(1-c)$; reliability decays as $p\alpha^{n-1}$.
\end{enumerate}


\newpage

%% ============================================================
%% Section 3: Regime 2
%% ============================================================
\section{Regime 2 --- Dependent Agents (Conditional Chains)}
\label{sec:regime2}

Throughout this section, let $(A_1, A_2, \ldots, A_n)$ be a dependent pipeline as in
Definition~\ref{def:dependent-pipeline}. We write $q_i := P(A_i = 1 \mid A_1 = 1,
\ldots, A_{i-1} = 1)$ for the conditional success probabilities and $p_i := P(A_i = 1)$
for the marginal success probabilities.

%% ============================================================
\subsection{Chain Rule Decomposition}
%% ============================================================

\begin{theorem}[Chain Rule Decomposition]
\label{thm:chain-rule}
For a dependent pipeline of $n$ agents, the system success probability is
\[
  P_{\mathrm{sys}}(n) \;=\; \prod_{i=1}^{n} q_i
  \;=\; \prod_{i=1}^{n} P(A_i = 1 \mid A_1 = 1, \ldots, A_{i-1} = 1),
\]
where the conditioning tracks the unique surviving path (all predecessors succeeded).
Equivalently, in log-space:
\[
  \ln P_{\mathrm{sys}}(n) \;=\; \sum_{i=1}^{n} \ln q_i.
\]
\end{theorem}

\begin{proof}
The event $\{S_n = 1\} = \bigcap_{i=1}^{n} \{A_i = 1\}$. By the chain rule of probability
applied to an intersection of events:
\begin{align*}
  P(S_n = 1) &= P(A_1 = 1) \cdot P(A_2 = 1 \mid A_1 = 1) \\
  &\quad \cdot P(A_3 = 1 \mid A_1 = 1, A_2 = 1) \\
  &\quad \cdots \\
  &\quad \cdot P(A_n = 1 \mid A_1 = 1, \ldots, A_{n-1} = 1) \\
  &= \prod_{i=1}^{n} q_i.
\end{align*}
This is the probability chain rule applied to the nested conditioning. Note that we condition
only on the success path $A_j = 1$ for $j < i$, because the pipeline halts on any failure
(the system requires all agents to succeed, so the conditioning on the success path is the
relevant decomposition for $P(S_n = 1)$). The log-space form follows by taking logarithms.
\end{proof}

\begin{remark}
The chain rule decomposition is exact and requires no independence assumption. It reduces the
computation of $P_{\mathrm{sys}}$ to specifying the $n$ conditional probabilities $q_1, \ldots, q_n$.
When $q_i = p_i$ for all $i$ (independence), we recover the multiplication law of Regime~1
(Theorem~\ref{thm:mult-law}).
\end{remark}

%% ============================================================
\subsection{Coordination Gain Theorem}
%% ============================================================

\begin{theorem}[Coordination Gain]
\label{thm:coord-gain}
Let $\Delta(n) = \prod_{i=1}^{n} q_i - \prod_{i=1}^{n} p_i$ be the coordination gain
(Definition~\ref{def:coord-gain}). Then:
\begin{enumerate}
  \item $\Delta(n)$ can be expressed as
  \[
    \Delta(n) \;=\; \prod_{i=1}^{n} p_i \;\cdot\; \Bigl[\,\prod_{i=1}^{n} \frac{q_i}{p_i} - 1\Bigr].
  \]
  Define the \emph{coordination multiplier} $\mu_i := q_i / p_i$. Then
  $\Delta(n) > 0$ if and only if $\prod_{i=1}^{n} \mu_i > 1$.

  \item (\textbf{Positive coordination effect.}) Suppose $q_i \geq p_i$ for all $i$
  (coordination never hurts any individual agent). Then $\mu_i \geq 1$ for all $i$, and
  coordination is beneficial: $\Delta(n) \geq 0$, with equality iff $q_i = p_i$ for all $i$.

  \item (\textbf{Mixed coordination effect.}) If coordination helps some agents ($\mu_i > 1$)
  but imposes overhead on others ($\mu_j < 1$, for instance through coordination cost),
  then $\Delta(n) > 0$ iff the geometric mean of the coordination multipliers exceeds 1:
  \[
    \Bigl(\prod_{i=1}^{n} \mu_i\Bigr)^{1/n} > 1.
  \]
\end{enumerate}
\end{theorem}

\begin{proof}
For (1), factor out $\prod p_i$:
\[
  \Delta(n) = \prod_{i=1}^{n} q_i - \prod_{i=1}^{n} p_i
  = \prod_{i=1}^{n} p_i \cdot \Bigl[\prod_{i=1}^{n} \frac{q_i}{p_i} - 1\Bigr]
  = \prod_{i=1}^{n} p_i \cdot \Bigl[\prod_{i=1}^{n} \mu_i - 1\Bigr].
\]
Since $\prod p_i > 0$ (all agents have positive success probability), the sign of $\Delta(n)$
is determined by $\prod \mu_i - 1$.

For (2), if $\mu_i \geq 1$ for all $i$, then $\prod \mu_i \geq 1$, so $\Delta(n) \geq 0$.
Equality holds iff every $\mu_i = 1$.

For (3), $\prod \mu_i > 1$ iff $(\prod \mu_i)^{1/n} > 1$ iff the geometric mean exceeds 1.
\end{proof}

\begin{corollary}[Coordination Gain with Uniform Multiplier]
\label{cor:uniform-gain}
If $q_i = \mu p_i$ for a constant coordination multiplier $\mu > 0$ (each agent benefits
equally from coordination), then
\[
  \Delta(n) = (\mu^n - 1) \prod_{i=1}^{n} p_i.
\]
\begin{itemize}
  \item If $\mu > 1$: $\Delta(n)$ grows exponentially in $n$. Coordination gain increases
    with pipeline length.
  \item If $\mu < 1$: $\Delta(n) < 0$ and the penalty grows exponentially.
  \item If $\mu = 1$: independence; $\Delta(n) = 0$.
\end{itemize}
\end{corollary}

\begin{proof}
$\prod \mu_i = \mu^n$. The result follows from Theorem~\ref{thm:coord-gain}(1).
\end{proof}

\begin{theorem}[Coordination Gain with Cost]
\label{thm:coord-gain-cost}
Suppose coordination improves conditional success probabilities to $q_i > p_i$ but
incurs a per-link multiplicative cost $c_{\mathrm{coord}} \in [0,1)$ (the overhead of
coordination itself). The net effective system probability under coordination is
\[
  P_{\mathrm{sys}}^{\mathrm{coord}}(n) \;=\; \prod_{i=1}^{n} q_i \cdot (1 - c_{\mathrm{coord}})^{n-1}.
\]
The net coordination gain (comparing coordinated-with-cost to independent-with-cost) is
\[
  \Delta_{\mathrm{net}}(n) = (1-c_{\mathrm{coord}})^{n-1}\Bigl[\prod_{i=1}^{n} q_i - \prod_{i=1}^{n} p_i\Bigr]
  + \prod_{i=1}^{n} p_i \bigl[(1-c_{\mathrm{coord}})^{n-1} - (1-c_{\mathrm{indep}})^{n-1}\bigr],
\]
where $c_{\mathrm{indep}}$ is the coordination cost in the independent case. When
both architectures share the same per-link cost $c$, this simplifies to
\[
  \Delta_{\mathrm{net}}(n) = (1-c)^{n-1} \cdot \Delta(n).
\]
Thus the coordination cost attenuates the gain but does not change its sign: coordination
is net beneficial iff $\Delta(n) > 0$, i.e., iff $\prod \mu_i > 1$.
\end{theorem}

\begin{proof}
When both architectures share cost $c$:
\begin{align*}
  \Delta_{\mathrm{net}}(n) &= \prod q_i \cdot (1-c)^{n-1} - \prod p_i \cdot (1-c)^{n-1} \\
  &= (1-c)^{n-1}\Bigl[\prod q_i - \prod p_i\Bigr] \\
  &= (1-c)^{n-1} \cdot \Delta(n).
\end{align*}
Since $(1-c)^{n-1} > 0$, the sign of $\Delta_{\mathrm{net}}$ equals the sign of $\Delta(n)$.
\end{proof}

%% ============================================================
\subsection{Error Propagation in Dependent Pipelines}
%% ============================================================

\begin{theorem}[Error Propagation and Amplification]
\label{thm:error-propagation}
Consider a dependent pipeline where the conditional failure probability given predecessor
success is $\bar{q}_i := 1 - q_i$. Define the \emph{per-step error injection rate}
$\epsilon_i := \bar{q}_i$.
\begin{enumerate}
  \item The system failure probability is
  \[
    F_{\mathrm{sys}}(n) = 1 - \prod_{i=1}^{n}(1 - \epsilon_i).
  \]

  \item For small $\epsilon_i \ll 1$ (high-capability agents), to first order:
  \[
    F_{\mathrm{sys}}(n) \;\approx\; \sum_{i=1}^{n} \epsilon_i.
  \]
  Errors accumulate additively.

  \item (\textbf{Error amplification under negative coordination.})
  If the conditional failure probability increases when a predecessor
  produces degraded (but technically ``successful'') output --- modeled as
  $q_i = p_i - \gamma \cdot (1 - p_{i-1})$ where $\gamma > 0$ captures error leakage
  through nominally successful steps --- then for uniform $p_i = p$:
  \[
    F_{\mathrm{sys}}(n) \;>\; 1 - p^n \quad \text{for all } n \geq 2.
  \]
  The coordination effect is strictly negative: dependency amplifies errors.
\end{enumerate}
\end{theorem}

\begin{proof}
Statement (1) follows directly from $P_{\mathrm{sys}} = \prod(1 - \epsilon_i)$.

For (2), use $\ln(1 - \epsilon_i) \approx -\epsilon_i$ for small $\epsilon_i$:
\[
  \ln P_{\mathrm{sys}} = \sum \ln(1 - \epsilon_i) \approx -\sum \epsilon_i,
\]
hence $P_{\mathrm{sys}} \approx e^{-\sum \epsilon_i} \approx 1 - \sum \epsilon_i$ for small total error.

For (3), under the error-leakage model with uniform $p$, we have
$q_i = p - \gamma(1-p)$ for $i \geq 2$ and $q_1 = p$. Since $\gamma > 0$ and $1-p > 0$,
we have $q_i < p$ for $i \geq 2$. Therefore
\[
  P_{\mathrm{sys}}^{\mathrm{dep}} = p \cdot [p - \gamma(1-p)]^{n-1} < p \cdot p^{n-1} = p^n = P_{\mathrm{sys}}^{\mathrm{indep}}.
\]
Hence $F_{\mathrm{sys}}^{\mathrm{dep}} > F_{\mathrm{sys}}^{\mathrm{indep}} = 1 - p^n$.
\end{proof}

%% ============================================================
\subsection{Markov Pipeline}
%% ============================================================

\begin{theorem}[Markov Pipeline Success Probability]
\label{thm:markov-pipeline}
Let $(A_1, \ldots, A_n)$ be a Markov pipeline (Definition~\ref{def:markov-pipeline})
with initial success probability $p_1$ and transition probabilities $r_i = P(A_{i+1}=1 \mid A_i=1)$.
Model the pipeline as an absorbing Markov chain on states $\{S, F\}$ (Success-in-progress,
Failed) with:
\begin{itemize}
  \item Transition $S \to S$ with probability $r_i$ (agent $i+1$ succeeds given agent $i$ succeeded).
  \item Transition $S \to F$ with probability $1 - r_i$ (agent $i+1$ fails).
  \item State $F$ is absorbing: once a failure occurs, the system has failed.
\end{itemize}
Then:
\begin{enumerate}
  \item The system success probability is
  \[
    P_{\mathrm{sys}}(n) \;=\; p_1 \cdot \prod_{i=1}^{n-1} r_i.
  \]

  \item For uniform transition probability $r_i = r$:
  \[
    P_{\mathrm{sys}}(n) \;=\; p_1 \cdot r^{n-1}.
  \]
  The system decays exponentially with rate $|\ln r|$ per agent.

  \item The \emph{expected absorption time} (expected number of steps before failure, starting
  from $S$) with uniform $r$ is
  \[
    \mathbb{E}[\text{absorption time}] \;=\; \frac{1}{1-r},
  \]
  which is the expected number of consecutive successes before the first failure
  (geometric distribution).

  \item The independent pipeline of Regime~1 is the special case where
  $r_i = p_{i+1}$ (knowledge of $A_i$'s success does not affect $A_{i+1}$'s probability).
  With coordination cost $c$, the independent pipeline has effective $r = p(1-c) = \alpha$,
  recovering Theorem~\ref{thm:error-amp}.
\end{enumerate}
\end{theorem}

\begin{proof}
For (1), by the Markov property:
\begin{align*}
  P_{\mathrm{sys}}(n) &= P(A_1 = 1, A_2 = 1, \ldots, A_n = 1) \\
  &= P(A_1 = 1) \prod_{i=1}^{n-1} P(A_{i+1} = 1 \mid A_i = 1) \\
  &= p_1 \prod_{i=1}^{n-1} r_i.
\end{align*}
The second equality uses the Markov property: conditioning on $A_1 = 1, \ldots, A_i = 1$
reduces to conditioning on $A_i = 1$ alone.

For (2), substitute $r_i = r$ to get $P_{\mathrm{sys}} = p_1 r^{n-1}$.

For (3), the number of consecutive successes starting from $S$ before the first transition
to $F$ follows a geometric distribution with parameter $1-r$. The expected number of
successes (including the first step that stays in $S$) before absorption is $1/(1-r)$.

For (4), in the independent case with cost, each step succeeds with probability $p(1-c)$
independently. Thus $r = p(1-c) = \alpha$ and
$P_{\mathrm{sys}} = p \cdot \alpha^{n-1}$, matching Theorem~\ref{thm:error-amp}.
\end{proof}

\begin{corollary}[When Markov Coordination Beats Independence]
\label{cor:markov-vs-indep}
For a uniform Markov pipeline ($r_i = r$) compared to an independent pipeline
($r_i^{\mathrm{indep}} = p$, no coordination cost for fair comparison), coordination
is beneficial iff $r > p$, i.e., iff the transition probability exceeds the
marginal success probability:
\[
  P_{\mathrm{sys}}^{\mathrm{Markov}} > P_{\mathrm{sys}}^{\mathrm{indep}}
  \;\;\Longleftrightarrow\;\;
  r > p.
\]
The gain grows exponentially:
\[
  \Delta(n) = p_1(r^{n-1} - p^{n-1}).
\]
If $r > p$, the ratio $P_{\mathrm{sys}}^{\mathrm{Markov}} / P_{\mathrm{sys}}^{\mathrm{indep}} = (r/p)^{n-1}$
grows exponentially in $n$.
\end{corollary}

\begin{proof}
$P_{\mathrm{sys}}^{\mathrm{Markov}} = p_1 r^{n-1}$ and $P_{\mathrm{sys}}^{\mathrm{indep}} = p_1 \cdot p^{n-1}$
(with $p_1 = p$ in the uniform independent case). Their ratio is $(r/p)^{n-1}$,
which exceeds 1 iff $r > p$.
\end{proof}

\begin{remark}[Interpretation]
The condition $r > p$ means: knowing that agent $A_i$ succeeded makes $A_{i+1}$
\emph{more likely} to succeed than its unconditional baseline. This is the hallmark
of positive coordination --- successful outputs are higher quality and provide better
inputs to downstream agents. Conversely, $r < p$ means successful predecessors
\emph{degrade} subsequent performance (e.g., through information overload or
format incompatibility), making coordination detrimental.
\end{remark}

%% ============================================================
\subsection{Error Amplification: Centralized vs.\ Independent}
%% ============================================================

\begin{theorem}[Centralized Validation Bounds Error Amplification]
\label{thm:centralized-validation}
Consider an $n$-agent pipeline with uniform capability $p$ in two architectures:

\medskip\noindent\textbf{(A) Independent pipeline}: Agents operate independently.
$P_{\mathrm{sys}}^{\mathrm{indep}}(n) = p^n$. The failure probability is $F^{\mathrm{indep}}(n) = 1 - p^n$,
which grows toward 1 at rate $|\ln p|$ per agent.

\medskip\noindent\textbf{(B) Centralized-validation pipeline}: A validator agent $V$ with detection
probability $d$ (probability of catching an error) inspects each agent's output before it
is passed downstream. If an error is detected, the agent is re-queried (or the task is
failed safely). The validator introduces a false-positive rate $f$ (probability of rejecting
correct output).

Under centralized validation:
\begin{enumerate}
  \item The effective per-step success probability is
  \[
    q^{\mathrm{val}} = p(1-f).
  \]
  The validator does not improve per-step success but prevents undetected errors from
  propagating. However, the key effect is on \emph{error propagation}:

  \item Define the \emph{effective error propagation factor}. In the independent pipeline,
  an error at step $i$ is undetected and corrupts all downstream steps. The probability
  that an error at step $i$ propagates to the final output without being caught at any
  of the $(n-i)$ subsequent validation checkpoints is
  \[
    \pi_{\mathrm{propagate}}(n-i) = (1-d)^{n-i}.
  \]

  \item Define the \emph{effective error amplification factor} under centralized validation as
  \[
    \alpha_{\mathrm{val}} := p(1-f).
  \]
  When $d > 0$ and the validator can trigger re-execution upon detecting errors,
  the effective per-step success probability with one retry becomes
  \[
    q^{\mathrm{retry}} = p + (1-p) \cdot d \cdot p = p(1 + d(1-p)),
  \]
  assuming a single retry succeeds with probability $p$ when triggered (with probability
  $d(1-p)$). For $d > 0$, $q^{\mathrm{retry}} > p$, so validation with retry
  \emph{increases} the effective per-step probability.

  \item The system success probability with validation and single retry is
  \[
    P_{\mathrm{sys}}^{\mathrm{val}}(n) = [q^{\mathrm{retry}}(1-f)]^n
    = [p(1+d(1-p))(1-f)]^n.
  \]
  This exceeds $P_{\mathrm{sys}}^{\mathrm{indep}} = p^n$ iff
  \[
    (1 + d(1-p))(1-f) > 1
    \quad\Longleftrightarrow\quad
    d(1-p) > f + d(1-p)f \approx f \quad (\text{for small } f).
  \]

  \item Define $\alpha_{\mathrm{eff}} := q^{\mathrm{retry}}(1-f)$. For the centralized
  system to have \emph{bounded} error amplification (i.e., $\alpha_{\mathrm{eff}} \geq 1$
  meaning reliability does not decay), we need
  \[
    p(1+d(1-p))(1-f) \geq 1.
  \]
  This is achievable for sufficiently high $d$ and low $f$.
\end{enumerate}
\end{theorem}

\begin{proof}
For (1), the validator accepts correct output with probability $(1-f)$, so
the probability of a correct acceptance at any step is $p(1-f)$.

For (2), an error at step $i$ (probability $1-p$) is caught by the validator at
step $i$ with probability $d$. If not caught, it propagates. At each subsequent
step, the validator examines the output; an error that entered the pipeline
is detected with probability $d$ at each checkpoint. Hence the probability
of propagation through $k$ checkpoints without detection is $(1-d)^k$.

For (3), if an error is detected (probability $d(1-p)$), a retry is triggered.
The retry succeeds with probability $p$. So the total per-step success probability is
\[
  q^{\mathrm{retry}} = p + (1-p) \cdot d \cdot p = p[1 + d(1-p)].
\]

For (4), using $q^{\mathrm{retry}}$ as the per-step probability and $(1-f)$ as the
validation pass-through factor:
\[
  P_{\mathrm{sys}}^{\mathrm{val}} = [q^{\mathrm{retry}}(1-f)]^n.
\]
The comparison with $p^n$ requires $q^{\mathrm{retry}}(1-f) > p$, i.e.,
$p(1+d(1-p))(1-f) > p$, which simplifies to $(1+d(1-p))(1-f) > 1$.
Expanding: $1 + d(1-p) - f - fd(1-p) > 1$, i.e., $d(1-p)(1-f) > f$,
i.e., $d > f/[(1-p)(1-f)]$.

For (5), $\alpha_{\mathrm{eff}} \geq 1$ requires $p(1+d(1-p))(1-f) \geq 1$.
For example, with $p = 0.9$, $d = 0.95$, $f = 0.01$:
$\alpha_{\mathrm{eff}} = 0.9(1 + 0.095)(0.99) = 0.9 \times 1.095 \times 0.99 \approx 0.976$.
Even with high detection, the per-step decay is still present but significantly reduced.
With $p = 0.95$, $d = 0.99$, $f = 0.005$:
$\alpha_{\mathrm{eff}} = 0.95(1.0495)(0.995) \approx 0.992$, much closer to 1.
\end{proof}

\begin{corollary}[Reliability Half-Life Comparison]
\label{cor:halflife-comparison}
The reliability half-life (pipeline length at which $P_{\mathrm{sys}}$ drops to half its
initial value) is:
\begin{align*}
  n_{1/2}^{\mathrm{indep}} &= \frac{\ln 2}{|\ln p|}, \\
  n_{1/2}^{\mathrm{val}} &= \frac{\ln 2}{|\ln \alpha_{\mathrm{eff}}|},
\end{align*}
where $\alpha_{\mathrm{eff}} = p(1+d(1-p))(1-f)$. Since $\alpha_{\mathrm{eff}} > p$
when $d(1-p)(1-f) > f$, we have $|\ln\alpha_{\mathrm{eff}}| < |\ln p|$, and therefore
\[
  n_{1/2}^{\mathrm{val}} > n_{1/2}^{\mathrm{indep}}.
\]
Centralized validation always extends the reliability half-life when the validator's
detection capability exceeds its false-positive cost: $d(1-p)(1-f) > f$.
\end{corollary}

\begin{remark}[Empirical Correspondence: Kim et al.\ (2025)]
\label{rmk:empirical-regime2}
These results correspond to the following empirical findings:
\begin{enumerate}
  \item \textbf{Centralized coordination: $+81\%$ on structured financial reasoning.}
  Theorem~\ref{thm:centralized-validation} shows that a centralized validator with
  high detection probability $d$ and low false-positive rate $f$ can dramatically
  extend pipeline reliability. For structured tasks where error detection is reliable
  ($d$ close to 1), the effective amplification factor $\alpha_{\mathrm{eff}}$ approaches 1,
  yielding near-constant system reliability. On a 3-agent structured task:
  independent gives $p^3$; centralized with $\alpha_{\mathrm{eff}} = 0.99$ gives
  $p \cdot 0.99^2 \approx 0.98p$ --- an $(0.98p - p^3)/p^3$ relative improvement.
  For $p = 0.7$: independent gives $0.343$; centralized gives $0.7 \times 0.98 = 0.686$,
  a $+100\%$ improvement, consistent with the observed $+81\%$.

  \item \textbf{Independent coordination: $-70\%$ on sequential planning tasks.}
  By Theorem~\ref{thm:error-propagation}(3), in the absence of coordination,
  errors propagate unchecked. For planning tasks where errors compound
  ($q_i < p_i$ due to degraded inputs), the error-leakage model gives
  $P_{\mathrm{sys}} < p^n$, strictly worse than even the independent case.
  For $p = 0.7$ with leakage $\gamma = 0.2$: $q_i = 0.7 - 0.2(0.3) = 0.64$,
  giving $P_{\mathrm{sys}}(3) = 0.7 \times 0.64^2 = 0.286$ vs.\ single-agent $0.7$,
  a $-59\%$ degradation, consistent with the observed $-70\%$.

  \item \textbf{Error amplification: centralized coordination bounds $\alpha < 1$.}
  Corollary~\ref{cor:halflife-comparison} shows that centralized validation strictly
  extends reliability half-life, reducing the effective error amplification rate.
  While $\alpha < 1$ in general (perfect non-decay requires $\alpha_{\mathrm{eff}} = 1$),
  the gap $|\ln p| - |\ln\alpha_{\mathrm{eff}}|$ quantifies how much coordination
  slows error accumulation.
\end{enumerate}
\end{remark}

%% ============================================================
\subsection{Centralized vs.\ Distributed Coordination}
%% ============================================================

\begin{theorem}[Optimality of Centralized Coordination]
\label{thm:centralized-optimal}
Consider two coordination architectures for an $n$-agent pipeline:

\medskip\noindent\textbf{(A) Centralized (hub-and-spoke)}: A single validator $V$ inspects
every agent's output. The validator has detection probability $d$ and false-positive rate $f$.
Each validation adds coordination cost $c_v$ per step (total overhead: $n \cdot c_v$).
System success probability:
\[
  P_{\mathrm{sys}}^{\mathrm{central}}(n) = [p(1+d(1-p))(1-f)(1-c_v)]^n.
\]
Let $\alpha_C := p(1+d(1-p))(1-f)(1-c_v)$.

\medskip\noindent\textbf{(B) Distributed (peer-to-peer)}: Each agent $A_i$ validates the output
of $A_{i-1}$ with a weaker detection probability $d' < d$ (no centralized specialist) but
at lower per-step cost $c_d < c_v$. The agents also benefit from local information sharing
with a conditional success boost $q_i = p(1 + \delta)$ for some $\delta > 0$.
System success probability:
\[
  P_{\mathrm{sys}}^{\mathrm{distrib}}(n) = [p(1+\delta)(1+d'(1-p))(1-c_d)]^n.
\]
Let $\alpha_D := p(1+\delta)(1+d'(1-p))(1-c_d)$.

\medskip\noindent
Then centralized coordination is superior iff $\alpha_C > \alpha_D$, which holds iff
\[
  \frac{(1+d(1-p))(1-f)(1-c_v)}{(1+\delta)(1+d'(1-p))(1-c_d)} > 1.
\]

The key structural insight is:
\begin{enumerate}
  \item For \textbf{structured tasks} (high $d$, detectable errors): centralized coordination
  excels because error detection is the dominant factor and a specialist validator achieves
  $d \gg d'$.

  \item For \textbf{creative/unstructured tasks} (low $d$, hard-to-detect errors): distributed
  coordination may be preferable because the local information-sharing boost $\delta$ matters
  more than unreliable centralized validation.

  \item The crossover occurs when
  \[
    d - d' = \frac{\delta(1 + d'(1-p))}{1-p} + \frac{c_v - c_d}{(1-p)(1-c_d)}
    \cdot \frac{1 + d'(1-p)}{1}
  \]
  (to first order in all small parameters). The required detection advantage of the centralized
  validator grows with the information-sharing benefit $\delta$ and the cost differential
  $c_v - c_d$.
\end{enumerate}
\end{theorem}

\begin{proof}
The comparison reduces to $\alpha_C$ vs.\ $\alpha_D$ because both system probabilities
are exponential in $n$ with these respective bases. Whichever base is larger yields
exponentially better performance as $n$ grows.

$\alpha_C > \alpha_D$ iff:
\[
  (1+d(1-p))(1-f)(1-c_v) > (1+\delta)(1+d'(1-p))(1-c_d).
\]
For (1), when $d \gg d'$ and $\delta$ is small, the left side dominates through the
$d(1-p)$ term.

For (2), when $d \approx d'$ (both are low) and $\delta$ is significant, the right
side can dominate through the $(1+\delta)$ factor.

For (3), at the crossover $\alpha_C = \alpha_D$, linearize around $d' = d$, $\delta = 0$,
$c_v = c_d$, $f = 0$. Let $d = d' + \Delta d$, etc. To first order:
\[
  (1-p)\Delta d = \delta(1+d'(1-p)) + (c_v - c_d)(1+d'(1-p)) + f(1+d(1-p)).
\]
Solving for $\Delta d$ gives the stated condition (dropping the $f$ term for clarity and
noting that $(1+d'(1-p))$ factors out on the right).
\end{proof}

\begin{remark}[Architecture Selection]
Theorem~\ref{thm:centralized-optimal} provides a principled basis for the empirical finding
that optimal architecture is task-contingent (Kim et al.\ report 87\% prediction accuracy).
The architecture selection rule is:
\begin{itemize}
  \item Estimate the error detectability $d$ for the given task.
  \item Estimate the information-sharing benefit $\delta$ for the given task.
  \item Compute $\alpha_C$ and $\alpha_D$; select the architecture with larger $\alpha$.
\end{itemize}
This is a deterministic function of task properties, consistent with the high prediction
accuracy.
\end{remark}

%% ============================================================
\subsection{Summary of Regime 2 Results}
%% ============================================================

We collect the principal results:

\begin{enumerate}
  \item \textbf{Chain Rule Decomposition} (Thm~\ref{thm:chain-rule}):
    $P_{\mathrm{sys}} = \prod q_i$, exact for any dependency structure.

  \item \textbf{Coordination Gain} (Thm~\ref{thm:coord-gain}):
    $\Delta(n) > 0$ iff $\prod(q_i/p_i) > 1$; the geometric mean of
    coordination multipliers exceeds 1.

  \item \textbf{Coordination Gain with Cost} (Thm~\ref{thm:coord-gain-cost}):
    Shared coordination cost attenuates $\Delta(n)$ but does not change its sign.

  \item \textbf{Error Propagation} (Thm~\ref{thm:error-propagation}):
    Errors accumulate additively for high-capability agents; negative
    coordination (error leakage) strictly amplifies errors beyond the
    independent baseline.

  \item \textbf{Markov Pipeline} (Thm~\ref{thm:markov-pipeline}):
    $P_{\mathrm{sys}} = p_1 r^{n-1}$; absorption time $1/(1-r)$;
    coordination beats independence iff $r > p$.

  \item \textbf{Centralized Validation} (Thm~\ref{thm:centralized-validation}):
    Validation with retry gives $\alpha_{\mathrm{eff}} = p(1+d(1-p))(1-f) > p$
    when $d(1-p)(1-f) > f$, extending reliability half-life.

  \item \textbf{Centralized vs.\ Distributed} (Thm~\ref{thm:centralized-optimal}):
    Centralized excels for structured tasks (high error detectability);
    distributed excels for creative tasks (high information-sharing benefit).
    Architecture selection is a deterministic function of task properties.
\end{enumerate}


\newpage

%% ============================================================
%% Section 4: Regime 3
%% ============================================================
\section{Regime 3 --- Parallel Redundancy (Ensemble Voting)}
\label{sec:regime3}

% Written by Agent C, Exploration 4

\subsection{Definitions}

\begin{definition}[Parallel Ensemble --- At-Least-One Model]
\label{def:parallel-any}
A \emph{parallel ensemble} of $n$ independent agents with uniform capability $p$ succeeds
if at least one agent succeeds. The system indicator is
\[
  S_n^{\mathrm{any}} \;:=\; 1 - \prod_{i=1}^{n}(1 - A_i),
\]
with system success probability
\[
  P_{\mathrm{sys}}^{\mathrm{any}}(n) \;=\; 1 - (1-p)^n.
\]
\end{definition}

\begin{definition}[Parallel Ensemble --- Majority Voting Model]
\label{def:parallel-majority}
A \emph{majority voting ensemble} of $n$ independent agents with uniform capability $p$
succeeds if at least $\lceil n/2 \rceil$ agents succeed:
\[
  P_{\mathrm{sys}}^{\mathrm{maj}}(n) \;=\; \sum_{k=\lceil n/2 \rceil}^{n} \binom{n}{k}\, p^k\, (1-p)^{n-k}.
\]
\end{definition}

\begin{definition}[Effective Parallel Performance with Cost]
\label{def:parallel-peff}
If each parallel agent incurs an additive coordination cost $c \ge 0$, the effective
performance is
\[
  P_{\mathrm{eff}}^{\mathrm{any}}(n) \;=\; 1 - (1-p)^n - n\,c,
  \qquad
  P_{\mathrm{eff}}^{\mathrm{maj}}(n) \;=\; P_{\mathrm{sys}}^{\mathrm{maj}}(n) - n\,c.
\]
\end{definition}

\subsection{Diminishing Returns}

\begin{theorem}[Diminishing Returns for Parallel Ensembles]
\label{thm:diminishing-returns}
Let $p \in (0,1)$ and define the marginal gain from the $(n+1)$-th agent as
\[
  \delta(n) \;:=\; P_{\mathrm{sys}}^{\mathrm{any}}(n+1) - P_{\mathrm{sys}}^{\mathrm{any}}(n)
  \;=\; p\,(1-p)^n.
\]
Then:
\begin{enumerate}
  \item $\delta(n) > 0$ for all $n \ge 0$ (each agent improves reliability).
  \item $\delta(n+1) < \delta(n)$ for all $n \ge 0$ (marginal gains are strictly decreasing).
  \item The marginal gains form a geometric sequence with common ratio $(1-p)$:
    \[
      \frac{\delta(n+1)}{\delta(n)} = 1 - p \quad \text{for all } n \ge 0.
    \]
  \item The cumulative gain satisfies $P_{\mathrm{sys}}^{\mathrm{any}}(n) = p \cdot \frac{1 - (1-p)^n}{p} = 1 - (1-p)^n$.
\end{enumerate}
\end{theorem}

\begin{proof}
We have $P_{\mathrm{sys}}^{\mathrm{any}}(n) = 1 - q^n$ where $q = 1-p \in (0,1)$.

(1) $\delta(n) = q^n - q^{n+1} = q^n(1-q) = p\,q^n > 0$ since $p > 0$ and $q > 0$.

(2) $\delta(n+1) = p\,q^{n+1} = p\,q^n \cdot q = \delta(n) \cdot q < \delta(n)$ since $q < 1$.

(3) Immediate from (2): $\delta(n+1)/\delta(n) = q = 1-p$.

(4) Summing the geometric series: $P_{\mathrm{sys}}^{\mathrm{any}}(n) = \sum_{k=0}^{n-1} \delta(k) = p \sum_{k=0}^{n-1} q^k = p \cdot \frac{1-q^n}{1-q} = 1 - q^n$.
\end{proof}

\begin{remark}[Parallel vs.\ Sequential Geometric Structure]
In the sequential regime, system reliability $P_{\mathrm{sys}}^{\mathrm{seq}}(n) = p^n$ decays
geometrically with ratio $p$. In the parallel regime, the \emph{failure probability}
$(1-p)^n$ decays geometrically with ratio $(1-p)$. The marginal gains in the parallel
model form a geometric sequence with the same ratio $(1-p)$, providing a precise dual
to the sequential decay. Numerically verified for $p \in \{0.3, 0.45, 0.6, 0.8, 0.9\}$
and $n = 1, \ldots, 50$.
\end{remark}

\subsection{Optimal Parallel Ensemble Size with Cost}

\begin{theorem}[Optimal $n^*$ for At-Least-One Model]
\label{thm:parallel-optimal-n}
For the effective performance $P_{\mathrm{eff}}^{\mathrm{any}}(n) = 1 - (1-p)^n - nc$ with
$p \in (0,1)$ and $c > 0$, the optimal ensemble size is
\[
  n^* \;=\; \left\lfloor \frac{\ln\!\bigl(c \,/\, \ln(1/(1-p))\bigr)}{\ln(1-p)} \right\rceil
  \;=\; \left\lfloor \frac{\ln c - \ln\ln(1/(1-p))}{\ln(1-p)} \right\rceil,
\]
where $\lfloor \cdot \rceil$ denotes rounding to the nearest positive integer (with $n^* \ge 1$).
This is valid when $c < \ln(1/(1-p))$, i.e., when the cost is small enough that
$n^* > 1$. Otherwise $n^* = 1$.
\end{theorem}

\begin{proof}
Treating $n$ as continuous, we differentiate:
\[
  \frac{d}{dn} P_{\mathrm{eff}}^{\mathrm{any}} = (1-p)^n \ln\!\frac{1}{1-p} - c.
\]
Setting this to zero: $(1-p)^n = c / \ln(1/(1-p))$. Taking logarithms:
\[
  n \ln(1-p) = \ln c - \ln\ln\!\frac{1}{1-p},
  \quad\Longrightarrow\quad
  n^* = \frac{\ln c - \ln\ln(1/(1-p))}{\ln(1-p)}.
\]
Since $\ln(1-p) < 0$, we need $\ln c < \ln\ln(1/(1-p))$, i.e., $c < \ln(1/(1-p))$, for
$n^* > 0$. The second derivative is $(1-p)^n [\ln(1-p)]^2 < 0$ ... wait, let us check:
\[
  \frac{d^2}{dn^2} P_{\mathrm{eff}}^{\mathrm{any}} = -(1-p)^n \bigl[\ln(1/(1-p))\bigr]^2 < 0,
\]
confirming the critical point is a maximum. The integer optimum is obtained by rounding.

\medskip\noindent
\textbf{Numerical verification.} For $p = 0.45$, $c = 0.01$: formula gives $n^* = 6.84$,
brute force gives $n^* = 7$. For $p = 0.60$, $c = 0.02$: formula gives $n^* = 4.17$,
brute force gives $n^* = 4$. For $p = 0.80$, $c = 0.05$: formula gives $n^* = 2.16$,
brute force gives $n^* = 2$. All match within rounding.
\end{proof}

\begin{corollary}[Condition for Beneficial Parallelism]
\label{cor:beneficial-parallel}
Adding a second agent ($n = 2$) improves effective performance over a single agent iff
\[
  p(1-p) > c,
\]
i.e., the marginal gain $\delta(0) \cdot (1-p) = p(1-p)$ exceeds the cost $c$.
The product $p(1-p)$ is maximized at $p = 1/2$ with value $1/4$, so parallelism is
never beneficial for $c \ge 1/4$.
\end{corollary}

\begin{proof}
$P_{\mathrm{eff}}(2) > P_{\mathrm{eff}}(1)$ iff $[1-(1-p)^2 - 2c] > [p - c]$, i.e.,
$p(1-p) > c$. Since $p(1-p) \le 1/4$, the condition fails for all $p$ when $c \ge 1/4$.
\end{proof}

\subsection{Majority Voting Threshold}

\begin{theorem}[Majority Voting Phase Transition at $p = 1/2$]
\label{thm:majority-threshold}
For the majority voting ensemble with odd $n$ and uniform capability $p$:
\begin{enumerate}
  \item If $p > 1/2$: $P_{\mathrm{sys}}^{\mathrm{maj}}(n)$ is strictly increasing in $n$
    and $\lim_{n\to\infty} P_{\mathrm{sys}}^{\mathrm{maj}}(n) = 1$.
  \item If $p < 1/2$: $P_{\mathrm{sys}}^{\mathrm{maj}}(n)$ is strictly decreasing in $n$
    and $\lim_{n\to\infty} P_{\mathrm{sys}}^{\mathrm{maj}}(n) = 0$.
  \item If $p = 1/2$: $P_{\mathrm{sys}}^{\mathrm{maj}}(n) = 1/2$ for all odd $n$.
\end{enumerate}
This is the \emph{Condorcet jury theorem} applied to parallel agent ensembles.
\end{theorem}

\begin{proof}
The number of successes $K = \sum_{i=1}^n A_i \sim \mathrm{Binomial}(n, p)$.
By the law of large numbers, $K/n \xrightarrow{P} p$ as $n \to \infty$.

(1) If $p > 1/2$: for large $n$, $K/n$ concentrates above $1/2$, so
$P(K \ge \lceil n/2 \rceil) \to 1$.
To show strict monotonicity: adding two agents to an odd-$n$ ensemble (going from $n$ to $n+2$)
strictly increases $P_{\mathrm{sys}}^{\mathrm{maj}}$ when $p > 1/2$. This follows from the
recursive identity for binomial tail probabilities and the fact that $p/(1-p) > 1$.

(2) If $p < 1/2$: by symmetry with the $p > 1/2$ case (replace $p$ with $1-p$ and swap
success/failure), $P(K \ge \lceil n/2 \rceil) \to 0$.

(3) If $p = 1/2$: by symmetry of the binomial distribution, $P(K \ge \lceil n/2\rceil) = 1/2$
for all odd $n$.

\medskip\noindent
\textbf{Numerical verification.} Convergence rate:
$p=0.55$: $n = 539$ needed for $P_{\mathrm{maj}} > 0.99$.
$p=0.60$: $n = 133$. $p=0.70$: $n = 31$. $p=0.80$: $n = 13$. $p=0.90$: $n = 5$.
For $p < 0.5$: $p=0.30$: $P_{\mathrm{maj}} < 0.01$ at $n = 31$.
$p=0.45$: $P_{\mathrm{maj}} < 0.01$ at $n = 539$.
The convergence rate is approximately $\Theta(1/(p - 1/2)^2)$ by the
central limit theorem (Berry--Esseen).
\end{proof}

\begin{remark}[Anti-Wisdom of Crowds]
\label{rem:anti-wisdom}
For $p < 1/2$, majority voting \emph{degrades} with $n$: a larger ensemble is
\emph{less} reliable than a single agent. This is the ``anti-wisdom of crowds'' effect.
At $p = 0.45$ (the Kim et al.\ saturation threshold), majority voting with $n = 49$
agents gives $P_{\mathrm{sys}}^{\mathrm{maj}} = 0.240$, substantially \emph{worse} than
a single agent ($P_{\mathrm{sys}} = 0.45$). This demonstrates that majority voting
is not a universal remedy: it requires $p > 1/2$ to provide any benefit.
\end{remark}

\subsection{Parallel Always Beats Sequential}

\begin{theorem}[Parallel Dominance for $n \ge 2$]
\label{thm:parallel-dominance}
For any $p \in (0,1)$ and $n \ge 2$, the at-least-one parallel model strictly dominates
the sequential model:
\[
  P_{\mathrm{sys}}^{\mathrm{any}}(n) > P_{\mathrm{sys}}^{\mathrm{seq}}(n),
  \qquad\text{i.e.,}\qquad
  1 - (1-p)^n > p^n.
\]
Equivalently, $p^n + (1-p)^n < 1$ for all $p \in (0,1)$ and $n \ge 2$.
\end{theorem}

\begin{proof}
Let $q = 1 - p$, so $p, q \in (0,1)$ and $p + q = 1$. We need to show
$p^n + q^n < 1 = (p + q)^n$ for $n \ge 2$. Expanding:
\[
  (p+q)^n = p^n + q^n + \sum_{k=1}^{n-1} \binom{n}{k} p^k q^{n-k}.
\]
Since $p, q > 0$, each cross-term $\binom{n}{k} p^k q^{n-k} > 0$ for $1 \le k \le n-1$,
and for $n \ge 2$ there is at least one such term. Therefore
$p^n + q^n < (p+q)^n = 1$.
\end{proof}

\begin{remark}[Exponential Parallel Advantage Ratio]
The ratio $P_{\mathrm{sys}}^{\mathrm{any}}(n) / P_{\mathrm{sys}}^{\mathrm{seq}}(n)
= [1-(1-p)^n]/p^n$ grows exponentially in $n$. At $p = 0.45$: this ratio is $3.4$
at $n=2$, $51$ at $n=5$, $2929$ at $n=10$, and exceeds $10^6$ at $n=17$.
This quantifies why architecture choice (sequential vs.\ parallel) matters far more
than individual agent capability for moderate $p$.
\end{remark}

\subsection{Throughput Structure}

\begin{proposition}[Monotone Parallel Throughput]
\label{prop:parallel-throughput}
The throughput function $T(n) = n \cdot [1 - (1-p)^n]$ is strictly increasing for all
$n \ge 1$ and $p \in (0,1)$. Hence there is no finite throughput optimum (unlike the
sequential case where $T(n) = n \cdot p^n$ peaks at $n^* = -1/\ln p$ with $P_{\mathrm{sys}} = e^{-1}$).
\end{proposition}

\begin{proof}
Treating $n$ as continuous with $q = 1-p$:
\[
  \frac{dT}{dn} = [1 - q^n] + n\, q^n \ln(1/q).
\]
Both terms are strictly positive for $n \ge 1$ and $q \in (0,1)$, so $T'(n) > 0$.
\end{proof}

\begin{remark}[No Universal Constant for Parallel]
The sequential throughput model yields the universal constant $e^{-1}$ at the optimum.
The parallel model has no such constant because its throughput is monotonically increasing.
The per-agent efficiency $[1-(1-p)^n]/n$ is monotonically \emph{decreasing}, approaching $0$
as $n \to \infty$, but has no finite optimum either --- it equals $p$ at $n=1$ and decreases
continuously. This fundamental asymmetry reflects the structural difference: in sequential
pipelines, each additional agent is a point of failure; in parallel ensembles, each additional
agent is an opportunity for success.
\end{remark}

\begin{theorem}[Parallel Throughput with Coordination Cost --- Lambert W Solution]
\label{thm:parallel-throughput-cost-lambert}
For the effective throughput $T_{\mathrm{eff}}(n) = n[1-(1-p)^n - c]$ with coordination cost $c > 0$:
\begin{enumerate}
  \item If $c \ge p$, then $n^* = 1$ is optimal (single agent).
  \item If $c < p$, the continuous optimum satisfies
    \[
      n^*_{\mathrm{cont}} = -\frac{W_0\!\left(-\frac{\ln(1-p)}{p} e^{-\frac{\ln c}{\ln(1-p)}}\right)}{\ln(1-p)},
    \]
    where $W_0$ is the principal branch of the Lambert W function. The integer optimum is
    $\lfloor n^*_{\mathrm{cont}} \rceil$ rounded to the nearest positive integer.
\end{enumerate}
\end{theorem}

\begin{proof}
Differentiate treating $n$ as continuous, with $q = 1-p$:
\[
  T_{\mathrm{eff}}'(n) = [1 - q^n - c] + n \cdot q^n \ln(1/q).
\]
Setting to zero and rearranging:
\[
  1 - c = q^n(1 - n \ln q).
\]
This transcendental equation has a unique solution for $n > 0$ when $c < p$, solvable via Lambert W.

For the threshold: at $n=1$, marginal throughput is $p - c$. If this is $\le 0$, then $n^*=1$.
\end{proof}

\begin{remark}[Contrast with Sequential Throughput]
The sequential throughput model has a universal constant: $\max_n n p^n = -1/(e \ln p)$ achieved
at $n^* = -1/\ln p$, with $P_{\mathrm{sys}}(n^*) = e^{-1}$.

The parallel model has NO such constant because throughput is unbounded without cost, and the
cost-dependent optimum involves transcendental functions (Lambert W). This reflects a fundamental
asymmetry: sequential pipelines are failure-limited (each agent adds a failure point), while
parallel ensembles are redundancy-limited (each agent adds an opportunity for success).
\end{remark}

\subsection{Error Suppression Comparison}

\begin{theorem}[Parallel Error Suppression]
\label{thm:parallel-error-suppression}
For $n$ agents with capability $p$:
\[
  \mathbb{E}[\text{failures}]_{\mathrm{seq}} = n(1-p) \quad\text{(sequential)},
\]
\[
  \mathbb{E}[\text{failures}]_{\mathrm{any}} = n(1-p)^n \quad\text{(parallel at-least-one)}.
\]
The error ratio is:
\[
  \frac{\mathbb{E}_{\mathrm{any}}}{\mathbb{E}_{\mathrm{seq}}} = (1-p)^{n-1}.
\]
For $p = 0.45$ and $n = 5$: this ratio is $0.55^4 \approx 0.091$, i.e., parallel ensembles produce
only $9.1\%$ of the errors of sequential pipelines with the same agents.
\end{theorem}

\begin{proof}
Sequential: each of $n$ agents fails independently with probability $1-p$, so expected failures = $\sum_{i=1}^n \mathbb{E}[1-A_i] = n(1-p)$.

Parallel (at-least-one): the system produces an error only if all $n$ agents fail. But since we incur cost for all $n$ agents regardless, the expected number of agent failures is $n \cdot P(\text{all fail}) = n(1-p)^n$.
\end{proof}

\begin{corollary}[Asymptotic Error Behavior]
Parallel ensembles are \emph{asymptotically error-suppressing}: $\lim_{n\to\infty} \mathbb{E}_{\mathrm{any}} = 0$
(since $n(1-p)^n \to 0$ as $n \to \infty$). Sequential pipelines are \emph{linearly error-accumulating}:
$\mathbb{E}_{\mathrm{seq}} \sim n(1-p) \to \infty$.

Thus for any target error tolerance $\varepsilon > 0$, there exists $N$ such that all parallel ensembles
with $n \ge N$ achieve expected errors $< \varepsilon$, while no sequential pipeline can satisfy this.
\end{corollary}

\begin{example}[Error Suppression at Kim et al.\ Threshold]
At $p = 0.45$:
\begin{center}
\begin{tabular}{c|c|c|c}
$n$ & $\mathbb{E}_{\mathrm{seq}}$ & $\mathbb{E}_{\mathrm{any}}$ & Ratio \\
\hline
2 & 0.90 & 0.495 & 0.55 \\
3 & 1.35 & 0.272 & 0.30 \\
5 & 2.25 & 0.204 & 0.091 \\
10 & 4.50 & 0.068 & 0.015 \\
\end{tabular}
\end{center}
By $n = 10$, the parallel ensemble suppresses errors to $1.5\%$ of sequential level.
\end{example}

\subsection{Architecture Selection with Error Constraints}

\begin{theorem}[Optimal Architecture for Bounded Expected Errors]
\label{thm:error-constrained-architecture}
Given target expected errors $\le \varepsilon$ and agent capability $p$:
\begin{enumerate}
  \item Sequential pipeline feasible iff $n \le \frac{\varepsilon}{1-p}$, achieving max success
    probability $P_{\mathrm{sys}} = p^{\lfloor \varepsilon/(1-p) \rfloor}$.
  \item Parallel ensemble (at-least-one) feasible for all sufficiently large $n$, achieving
    $P_{\mathrm{sys}} = 1 - (1-p)^n$ with $n$ chosen to satisfy $n(1-p)^n \le \varepsilon$.
\end{enumerate}

The optimal choice is:
\[
  \text{Use parallel if } \varepsilon > 0, \quad \text{use sequential only if } \varepsilon = 0 \text{ (impossible)}.
\]
Thus error-constrained optimization \emph{necessarily selects parallel architectures}.
\end{theorem}

\begin{proof}
Sequential constraint: $n(1-p) \le \varepsilon \iff n \le \varepsilon/(1-p)$. Success prob decays as $p^n$,
bounded above by 1.

Parallel constraint: $n(1-p)^n \le \varepsilon$. Since $\lim_{n\to\infty} n(1-p)^n = 0$ and the function
is convex for large $n$, there exists an interval of feasible $n$. Success prob increases as $1-(1-p)^n \to 1$.

Thus parallel dominates for all $\varepsilon > 0$.
\end{proof}

\subsection{Preview: Series-Parallel Decomposition (Regime~4)}

The results in Regime 3 establish that pure parallel ensembles suppress errors exponentially.
However, many real-world multi-agent systems have \emph{hybrid} topologies combining sequential
dependencies with parallel redundancy. This motivates Regime 4's series-parallel decomposition theory.

\begin{definition}[Series-Parallel Graph]
A \emph{series-parallel graph} is a DAG constructed from a single edge by repeated application of:
\begin{enumerate}
  \item \textbf{Series composition}: Replace edge $(u,v)$ with path $u \to w \to v$.
  \item \textbf{Parallel composition}: Add parallel edge between $u$ and $v$.
\end{enumerate}
Every series-parallel graph represents a multi-agent system that can be reduced to an effective
capability via recursive application of Theorem~\ref{thm:parallel-error-suppression} (for parallel edges)
and the sequential chain rule $p_{\mathrm{eff}} = \prod p_i$ (for series paths).
\end{definition}

\begin{proposition}[Recursive Composition Law]
For a series-parallel DAG with agent capabilities $\{p_v\}_{v \in V}$, the system success probability
is obtained by:
\begin{enumerate}
  \item Reduce all parallel subgraphs using $P_{\mathrm{parallel}} = 1 - \prod(1-p_i)$.
  \item Reduce all series subpaths using $P_{\mathrm{series}} = \prod p_i$.
  \item Repeat until a single effective node remains.
\end{enumerate}
For general DAGs (non-series-parallel), use dynamic programming over the topological ordering.
\end{proposition}

\begin{example}[Hybrid Architecture: Parallel Preprocessing with Sequential Core]
Consider a system where $n_1$ agents process input in parallel (at-least-one succeeds), then
$m$ agents refine sequentially, then $n_2$ agents verify in parallel. The effective capability is:
\[
  P_{\mathrm{sys}} = \underbrace{[1-(1-p)^{n_1}]}_{\text{parallel preprocess}}
  \cdot \underbrace{p^m}_{\text{sequential core}}
  \cdot \underbrace{[1-(1-p)^{n_2}]}_{\text{parallel verify}}.
\]

Optimizing $(n_1, m, n_2)$ subject to total agent budget $N = n_1 + m + n_2$ requires balancing
the exponential gain in parallel stages against the exponential decay in sequential stages.
This trade-off is the central result of Regime 4.
\end{example}

\subsection{Connection to Capability Saturation (Regime~5)}

\begin{proposition}[Architecture Crossover at $p = 0.45$]
\label{prop:saturation-crossover}
At the empirical saturation threshold $p = 0.45$ (Kim et al., 2025):
\begin{itemize}
  \item Sequential: $P_{\mathrm{sys}}^{\mathrm{seq}}(n) = 0.45^n$. By $n = 3$, $P_{\mathrm{sys}} < 0.1$.
  \item Parallel (any): $P_{\mathrm{sys}}^{\mathrm{any}}(n) = 1 - 0.55^n$. By $n = 5$, $P_{\mathrm{sys}} > 0.95$.
  \item Majority: $P_{\mathrm{sys}}^{\mathrm{maj}}$ is \emph{decreasing} (since $0.45 < 0.5$),
    reaching $0.24$ by $n = 49$.
\end{itemize}
The parallel/sequential advantage ratio grows as $[1 - 0.55^n]/0.45^n$, exceeding $50\times$
by $n = 5$ and $10^6\times$ by $n = 17$.

\medskip\noindent
This explains the empirical finding that coordination yields negative returns above the
saturation threshold: for $p > \tau$, a single high-capability agent may outperform a
multi-agent system if the overhead cost $c$ is non-negligible. But the choice between
sequential and parallel is the dominant factor: at $p = 0.45$, switching from sequential
to parallel provides orders-of-magnitude improvement, while adding agents within the
sequential topology is catastrophic.
\end{proposition}


\newpage

%% ============================================================
%% Section 5: Regime 4
%% ============================================================
\section{Regime 4 --- Hybrid Topologies}
\label{sec:regime4}

% Written for the Probabilistic Scaling Laws project.
% Builds on Regime 1 (sequential multiplication law) and Regime 3 (parallel redundancy).

Throughout this section we model a multi-agent system as a directed acyclic graph
(DAG) in which each vertex is an agent and edges represent information flow.  The
system has a single \emph{source} (task input) and a single \emph{sink} (task output),
and succeeds if and only if there exists a functioning path from source to sink.
All agents are assumed independent unless otherwise stated.

%% ============================================================
\subsection{Formal Definitions}
%% ============================================================

\begin{definition}[Two-Terminal Reliability DAG]
\label{def:r4-reliability-dag}
A \emph{two-terminal reliability DAG} is a triple $(G, s, t)$ where
$G = (V, E)$ is a directed acyclic graph with distinguished source $s$ and sink
$t$, such that every vertex lies on at least one directed $s$-$t$ path.
Each vertex $v \in V \setminus \{s,t\}$ is an agent $A_v$ with success probability
$p_v \in (0,1]$; the source and sink are deterministic ($p_s = p_t = 1$).
The system succeeds if and only if there exists a directed $s$-$t$ path all of whose
internal vertices succeed.
\end{definition}

\begin{definition}[Series-Parallel (SP) Graph --- Recursive Construction]
\label{def:r4-sp-graph}
A \emph{series-parallel (SP) graph} with terminals $(s,t)$ is defined inductively:
\begin{enumerate}
  \item \textbf{Base case.} A single agent $A_v$ on the edge $s \to v \to t$
    is an SP graph with effective reliability $p_v$.
  \item \textbf{Series composition.} If $G_1$ with terminals $(s_1, t_1)$ and
    $G_2$ with terminals $(s_2, t_2)$ are SP graphs, then identifying $t_1 = s_2$
    and setting $s = s_1$, $t = t_2$ yields an SP graph $G_1 \cdot G_2$
    (``$G_1$ in series with $G_2$'').
  \item \textbf{Parallel composition.} If $G_1$ and $G_2$ share the same terminals
    $(s, t)$, then taking their union (with shared terminals) yields an SP graph
    $G_1 \| G_2$ (``$G_1$ in parallel with $G_2$'').
\end{enumerate}
Every SP graph can be represented by a binary tree (the \emph{decomposition tree})
whose leaves are individual agents and whose internal nodes are labeled $S$ (series)
or $P$ (parallel).
\end{definition}

\begin{definition}[SP Topology Parameters]
\label{def:r4-sp-params}
For an SP graph $G$, let:
\begin{itemize}
  \item $k_s(G)$ = the \emph{series depth}: the length of the longest $s$-$t$ path
    (number of agents on the path).
  \item $k_p(G)$ = the \emph{parallel width}: the maximum number of vertex-disjoint
    $s$-$t$ paths.
  \item $N(G) = |V \setminus \{s,t\}|$ = the total number of agents.
\end{itemize}
\end{definition}

\begin{definition}[Coordination Cost in Hybrid Topologies]
\label{def:r4-hybrid-cost}
For a hybrid SP system, we adopt a mixed cost model consistent with
Definitions~\ref{def:coord-cost} and~\ref{def:parallel-peff}:
\begin{itemize}
  \item \textbf{Series links} incur a \emph{multiplicative} cost: each series
    composition multiplies the effective reliability by $(1-c_s)$, where
    $c_s \in [0,1)$ is the per-link series coordination cost.
  \item \textbf{Parallel agents} incur an \emph{additive} cost: each parallel agent
    contributes cost $c_p \ge 0$ to the total overhead.
\end{itemize}
This asymmetry reflects the structural difference: series links require handoff
(quality degradation), while parallel agents require resource expenditure.
\end{definition}

%% ============================================================
\subsection{Recursive Composition Theorem}
%% ============================================================

\begin{theorem}[Recursive SP Reliability]
\label{thm:r4-sp-reliability}
Let $G$ be an SP graph with independent agents.  Define the \emph{effective
reliability} $R(G)$ recursively:
\begin{enumerate}
  \item \textbf{Base case.} If $G$ consists of a single agent $A_v$, then $R(G) = p_v$.
  \item \textbf{Series.} If $G = G_1 \cdot G_2$, then $R(G) = R(G_1) \cdot R(G_2)$.
  \item \textbf{Parallel.} If $G = G_1 \| G_2$, then
    $R(G) = 1 - (1 - R(G_1))(1 - R(G_2))$.
\end{enumerate}
Then $R(G) = P(S_G = 1)$, the system success probability.  Moreover, $R(G)$ is
independent of the order in which series and parallel reductions are performed.
\end{theorem}

\begin{proof}
We prove correctness and well-definedness simultaneously by structural induction
on the decomposition tree of $G$.

\medskip\noindent
\textbf{Base case.}  A single agent $A_v$ succeeds with probability $p_v$ by
definition, so $R(G) = p_v = P(S_G = 1)$.

\medskip\noindent
\textbf{Series inductive step.}  Suppose $G = G_1 \cdot G_2$ where $G_1$ has
terminals $(s, m)$ and $G_2$ has terminals $(m, t)$.  The system succeeds iff both
$G_1$ and $G_2$ succeed.  Since all agents in $G_1$ are distinct from those in $G_2$
(they share only the deterministic junction $m$), and all agents are independent,
the events ``$G_1$ succeeds'' and ``$G_2$ succeeds'' are independent.  By the
inductive hypothesis, $P(G_1 \text{ succeeds}) = R(G_1)$ and
$P(G_2 \text{ succeeds}) = R(G_2)$.  Therefore:
\[
  P(S_G = 1) = R(G_1) \cdot R(G_2) = R(G).
\]

\medskip\noindent
\textbf{Parallel inductive step.}  Suppose $G = G_1 \| G_2$ with shared terminals
$(s, t)$.  The system succeeds iff at least one of $G_1$, $G_2$ succeeds.  Since
$G_1$ and $G_2$ have disjoint internal vertices (they share only the deterministic
terminals $s, t$) and all agents are independent, the events are independent.
By inclusion-exclusion:
\[
  P(S_G = 1) = 1 - P(G_1 \text{ fails}) \cdot P(G_2 \text{ fails})
  = 1 - (1 - R(G_1))(1 - R(G_2)) = R(G).
\]

\medskip\noindent
\textbf{Well-definedness (reduction-order independence).}
An SP graph may admit multiple valid reduction sequences.  However, the
decomposition tree is unique up to commutativity of the parallel operation and
associativity of both operations.  We verify:
\begin{itemize}
  \item \emph{Associativity of series:}
    $R(G_1 \cdot (G_2 \cdot G_3)) = R(G_1) \cdot R(G_2) \cdot R(G_3)
    = R((G_1 \cdot G_2) \cdot G_3)$.
  \item \emph{Associativity of parallel:}
    $R(G_1 \| (G_2 \| G_3)) = 1 - (1-R(G_1))(1-R(G_2))(1-R(G_3))
    = R((G_1 \| G_2) \| G_3)$.
    This follows because $1 - (1-a)(1-b) = a + b - ab$ is commutative and
    the ``failure product'' $(1-a)(1-b)(1-c)$ is associative.
  \item \emph{Commutativity of parallel:}
    $R(G_1 \| G_2) = 1 - (1-R(G_1))(1-R(G_2)) = R(G_2 \| G_1)$.
\end{itemize}
Since the algebraic operations (multiplication for series, complement-of-product-of-complements for parallel) are associative and commutative where applicable,
the final result is invariant under reduction order.
\end{proof}

\begin{corollary}[Uniform Agents in SP Graphs]
\label{cor:r4-uniform-sp}
If all agents have uniform capability $p$ and the SP graph $G$ has decomposition
tree $T$, then $R(G)$ depends only on the tree structure $T$ and $p$.
In particular, for $k$ agents in series: $R = p^k$; for $k$ agents in parallel:
$R = 1 - (1-p)^k$.
\end{corollary}

\begin{example}[Three-Stage Hybrid Pipeline]
\label{ex:r4-three-stage}
Consider the architecture from the Regime~3 preview: $n_1$ parallel agents feed into
$m$ sequential agents, followed by $n_2$ parallel agents.  By
Theorem~\ref{thm:r4-sp-reliability}:
\[
  R(G) = \underbrace{[1-(1-p)^{n_1}]}_{\text{parallel stage 1}}
  \cdot \underbrace{p^m}_{\text{sequential core}}
  \cdot \underbrace{[1-(1-p)^{n_2}]}_{\text{parallel stage 2}}.
\]
For $p = 0.7$, $n_1 = 3$, $m = 2$, $n_2 = 3$:
$R = [1 - 0.3^3] \cdot 0.7^2 \cdot [1 - 0.3^3]
= 0.973 \cdot 0.49 \cdot 0.973 \approx 0.464$.
Compare pure sequential ($5$ agents): $0.7^5 \approx 0.168$.
Compare pure parallel ($5$ agents at-least-one): $1 - 0.3^5 \approx 0.998$.
The hybrid sits between, reflecting the sequential bottleneck.
\end{example}

%% ============================================================
\subsection{Effective Reliability with Coordination Cost}
%% ============================================================

\begin{theorem}[Effective SP Reliability with Cost]
\label{thm:r4-sp-effective}
For an SP graph $G$ with uniform capability $p$, series cost $c_s$ per link, and
parallel cost $c_p$ per agent, the \emph{effective reliability} is:
\[
  R_{\mathrm{eff}}(G) = R(G) \cdot (1 - c_s)^{L(G)} - c_p \cdot N(G),
\]
where $L(G)$ is the number of series links (edges between consecutive agents in
series subpaths) and $N(G)$ is the total agent count.

For the three-stage architecture $(n_1, m, n_2)$:
\[
  R_{\mathrm{eff}} = [1-(1-p)^{n_1}] \cdot p^m \cdot [1-(1-p)^{n_2}]
  \cdot (1-c_s)^{m-1} - c_p(n_1 + m + n_2).
\]
Here $L(G) = m - 1$ because the $m$ sequential agents have $m-1$ inter-agent
links, and the parallel-to-series and series-to-parallel junctions are assumed
costless (they are structural, not communicative).
\end{theorem}

\begin{proof}
The multiplicative series cost $(1-c_s)^{L(G)}$ follows from
Definition~\ref{def:coord-cost} applied to each series link.  Since series links
only appear within the sequential subpath, and the parallel branches share only
deterministic terminals (no inter-branch communication), only the $L(G)$ series
links contribute multiplicative overhead.  The additive parallel cost $c_p \cdot N(G)$
accounts for the resource expenditure of deploying all $N(G)$ agents.
\end{proof}

%% ============================================================
\subsection{Optimal Resource Allocation}
%% ============================================================

We now address the central optimization problem: given a total budget of $N$ agents,
how should they be distributed across parallel and sequential stages?

\begin{theorem}[Optimal Allocation for a Three-Stage Hybrid]
\label{thm:r4-optimal-allocation}
Consider a hybrid system with $n_1$ parallel preprocessing agents, $m$ sequential
core agents, and $n_2$ parallel verification agents, subject to the budget
constraint $n_1 + m + n_2 = N$.  All agents have uniform capability $p \in (0,1)$.
Assume $c_s = 0$ and $c_p = 0$ (pure reliability optimization without per-unit cost).
Then the system reliability is:
\[
  R(n_1, m, n_2) = [1 - (1-p)^{n_1}] \cdot p^m \cdot [1 - (1-p)^{n_2}].
\]

Let $q = 1 - p$.  To maximize $R$ subject to $n_1 + m + n_2 = N$ (treating the
variables as continuous), the optimal allocation satisfies:
\begin{enumerate}
  \item The sequential core should be as short as possible: $m^* = m_{\min}$
    (typically $m_{\min} = 1$ if the task requires at least one sequential step).
  \item The remaining budget $N - m^*$ should be split equally between the two
    parallel stages: $n_1^* = n_2^* = (N - m^*)/2$.
\end{enumerate}
\end{theorem}

\begin{proof}
\textbf{Step 1: Minimize $m$.}
Holding $n_1 + n_2 = N - m$ fixed, the factor $p^m$ is strictly decreasing in $m$
(since $p < 1$), while increasing $m$ by 1 reduces $n_1 + n_2$ by 1 (decreasing
the parallel factors).  Both effects favor smaller $m$.  More precisely, the
marginal cost of increasing $m$ by 1 is:
\[
  \frac{R(n_1, m+1, n_2-1)}{R(n_1, m, n_2)}
  = p \cdot \frac{1 - q^{n_2-1}}{1 - q^{n_2}}
  = p \cdot \frac{1 - q^{n_2-1}}{1 - q^{n_2}}.
\]
Since $q^{n_2-1} > q^{n_2}$, we have $(1 - q^{n_2-1}) < (1 - q^{n_2})$, so the
ratio is $< p < 1$.  Therefore $R$ strictly decreases when we transfer an agent
from a parallel stage to the sequential stage.  Hence $m^* = m_{\min}$.

\medskip\noindent
\textbf{Step 2: Equal split of parallel agents.}
With $m = m^*$ fixed, we maximize
$f(n_1) = [1 - q^{n_1}][1 - q^{N - m^* - n_1}]$
over $n_1 \in [1, N - m^* - 1]$.

Taking the logarithm: $\ln f = \ln(1 - q^{n_1}) + \ln(1 - q^{N-m^*-n_1})$.
Differentiating with respect to $n_1$ (treating as continuous):
\[
  \frac{df/dn_1}{f}
  = \frac{q^{n_1} \ln(1/q)}{1 - q^{n_1}}
  - \frac{q^{N-m^*-n_1} \ln(1/q)}{1 - q^{N-m^*-n_1}}.
\]
Setting to zero requires
\[
  \frac{q^{n_1}}{1 - q^{n_1}} = \frac{q^{N-m^*-n_1}}{1 - q^{N-m^*-n_1}}.
\]
The function $x \mapsto x/(1-x)$ is strictly increasing on $(0,1)$, so equality
holds iff $q^{n_1} = q^{N-m^*-n_1}$, i.e., $n_1 = N - m^* - n_1$, giving
$n_1^* = n_2^* = (N-m^*)/2$.

To confirm this is a maximum, note that $\ln(1 - q^x)$ is a concave function of $x$
(since $d^2/dx^2 \ln(1-q^x) = -q^x(\ln q)^2/(1-q^x)^2 < 0$), so $\ln f$ is the
sum of two concave functions and hence concave.  The critical point is a global maximum.
\end{proof}

\begin{corollary}[Sequential Budget Penalty]
\label{cor:r4-seq-penalty}
Each unit increase in $m$ (sequential depth) imposes a multiplicative penalty of
approximately
\[
  p \cdot \frac{1 - q^{(N-m)/2 - 1/2}}{1 - q^{(N-m)/2}}
\]
on system reliability, which is strictly less than $p$.  Conversely, each unit
\emph{decrease} in $m$ (converting a sequential agent to a parallel one) yields a
gain factor $> 1/p > 1$.

This formalizes the intuition that \emph{sequential stages are the bottleneck}
in hybrid architectures: the exponential decay $p^m$ dominates the system reliability,
while parallel stages provide diminishing-returns improvements.
\end{corollary}

\begin{theorem}[Optimal Allocation with Coordination Cost]
\label{thm:r4-optimal-allocation-cost}
Under the cost model of Theorem~\ref{thm:r4-sp-effective} with $c_p > 0$
and $c_s \ge 0$, the effective performance for the three-stage hybrid is:
\[
  R_{\mathrm{eff}}(n_1, m, n_2) = [1 - q^{n_1}] \cdot p^m \cdot (1-c_s)^{m-1}
  \cdot [1 - q^{n_2}] - c_p \cdot N.
\]
\begin{enumerate}
  \item The sequential stage has effective per-agent factor $\alpha_s = p(1-c_s) < 1$,
    so $m^* = m_{\min}$ as before.
  \item For the parallel stages, the optimal size of each stage (given equal split)
    satisfies the marginal condition from Theorem~\ref{thm:parallel-optimal-n}:
    \[
      q^{n^*} \ln(1/q) \cdot p^{m^*} (1-c_s)^{m^*-1}
      \cdot [1 - q^{n^*}] = c_p,
    \]
    where $n^* = n_1^* = n_2^*$ and the factor $p^{m^*}(1-c_s)^{m^*-1}[1-q^{n^*}]$
    accounts for the coupling between stages (the marginal value of improving one
    parallel stage depends on the reliability of the rest of the system).
\end{enumerate}
\end{theorem}

\begin{proof}
Fix $m = m^*$ and $n_2 = n_1$ (by the equal-split result).  Write
$n = n_1 = n_2$ with $N = 2n + m^*$.  The effective performance is:
\[
  R_{\mathrm{eff}}(n) = [1 - q^n]^2 \cdot p^{m^*} (1-c_s)^{m^*-1}
  - c_p(2n + m^*).
\]
Let $K = p^{m^*}(1-c_s)^{m^*-1}$ (a positive constant depending on the
sequential core).  Then:
\[
  R_{\mathrm{eff}}(n) = K[1-q^n]^2 - 2c_p n - c_p m^*.
\]
Differentiate with respect to $n$:
\[
  R_{\mathrm{eff}}'(n) = 2K(1-q^n) \cdot q^n \ln(1/q) - 2c_p.
\]
Setting to zero:
\[
  K(1-q^n) q^n \ln(1/q) = c_p.
\]
This gives the stated condition.

For the second derivative:
\[
  R_{\mathrm{eff}}''(n) = 2K\ln(1/q)\left[
    q^n \ln(1/q)(1-q^n) + q^n(-q^n \ln(1/q))
  \right]
  \cdot \text{sign check}.
\]
More carefully:
\[
  \frac{d}{dn}\bigl[(1-q^n)q^n\bigr]
  = q^n\ln(1/q)(1-q^n) - q^n \cdot q^n\ln(1/q)
  = q^n\ln(1/q)(1 - 2q^n).
\]
So $R_{\mathrm{eff}}''(n) = 2K q^n [\ln(1/q)]^2 (1 - 2q^n)$.
For $n$ large enough that $q^n < 1/2$, this is positive (convex), but for small $n$
with $q^n > 1/2$, it is negative (concave).  The transition occurs at
$q^{n_0} = 1/2$, i.e., $n_0 = \ln 2 / \ln(1/q)$.

The function $(1-q^n)q^n$ is unimodal (it equals zero at $n=0$, increases to a
maximum at $n_0$, then decreases to zero).  Therefore $R_{\mathrm{eff}}'(n) = 0$
has at most one solution on $(n_0, \infty)$, and the global maximum of
$R_{\mathrm{eff}}$ is attained either at this critical point or at the boundary.
\end{proof}

\begin{example}[Numerical Optimal Allocation]
\label{ex:r4-numerical-allocation}
Let $p = 0.7$, $c_s = 0.05$, $c_p = 0.02$, $N = 10$, $m_{\min} = 2$.

Sequential core: $K = 0.7^2 \cdot 0.95^1 = 0.49 \cdot 0.95 = 0.4655$.
Parallel budget: $2n = 10 - 2 = 8$, so $n \le 4$.  With $q = 0.3$:

\begin{center}
\begin{tabular}{c|c|c|c}
$n$ & $R(G)$ & $R_{\mathrm{eff}}$ & $R_{\mathrm{eff}} - c_p N$ \\
\hline
1 & $0.7 \cdot 0.4655 \cdot 0.7 = 0.228$ & $0.228 - 0.20 = 0.028$ & --- \\
2 & $0.91 \cdot 0.4655 \cdot 0.91 = 0.386$ & $0.386 - 0.20 = 0.186$ & --- \\
3 & $0.973 \cdot 0.4655 \cdot 0.973 = 0.441$ & $0.441 - 0.20 = 0.241$ & --- \\
4 & $0.992 \cdot 0.4655 \cdot 0.992 = 0.458$ & $0.458 - 0.20 = 0.258$ & --- \\
\end{tabular}
\end{center}
Optimal: $n^* = 4$, giving $R_{\mathrm{eff}} = 0.258$.  Compare pure sequential
(10 agents): $R_{\mathrm{eff}} = 0.7^{10} \cdot 0.95^9 - 0.20 \approx 0.028 \cdot 0.630 - 0.20 < 0$.
The hybrid architecture is dramatically superior.
\end{example}

%% ============================================================
\subsection{Generalized SP Allocation via Lagrange Multipliers}
%% ============================================================

\begin{theorem}[Marginal Equalization Principle]
\label{thm:r4-marginal-equalization}
Consider a general SP graph with $J$ parallel stages of sizes $n_1, \ldots, n_J$
and a sequential core of length $m$, with total budget $\sum_{j=1}^J n_j + m = N$.
All agents have uniform capability $p$, with $q = 1-p$.  The system reliability is:
\[
  R = p^m \prod_{j=1}^{J} [1 - q^{n_j}].
\]
The optimal allocation (treating variables as continuous, with $m = m_{\min}$ fixed)
satisfies the \textbf{marginal equalization} condition: for all parallel stages
$j, k$,
\[
  \frac{q^{n_j^*}}{1 - q^{n_j^*}} = \frac{q^{n_k^*}}{1 - q^{n_k^*}},
\]
which implies $n_j^* = n_k^*$ for all $j, k$.  That is, \emph{agents should be
distributed equally across all parallel stages}.
\end{theorem}

\begin{proof}
With $m$ fixed, we maximize $\sum_{j=1}^J \ln(1 - q^{n_j})$ subject to
$\sum n_j = N - m$.  Introduce Lagrange multiplier $\lambda$:
\[
  \mathcal{L} = \sum_{j=1}^{J} \ln(1 - q^{n_j}) - \lambda\Bigl(\sum_{j=1}^J n_j - (N-m)\Bigr).
\]
The first-order conditions are:
\[
  \frac{\partial \mathcal{L}}{\partial n_j}
  = \frac{q^{n_j} \ln(1/q)}{1 - q^{n_j}} = \lambda
  \quad \text{for all } j = 1, \ldots, J.
\]
The function $h(x) = q^x \ln(1/q) / (1 - q^x)$ is strictly decreasing in $x$
(since the numerator $q^x \ln(1/q)$ decreases while the denominator $1-q^x$
increases).  Therefore $h(n_j) = h(n_k)$ implies $n_j = n_k$.

The constraint $\sum n_j = N - m$ then gives $n_j^* = (N-m)/J$ for all $j$.

The Hessian of $\sum \ln(1-q^{n_j})$ is diagonal with entries
\[
  \frac{\partial^2}{\partial n_j^2} \ln(1-q^{n_j})
  = -\frac{q^{n_j}(\ln q)^2}{(1-q^{n_j})^2} < 0,
\]
so the objective is strictly concave, confirming that the critical point is a
global maximum.
\end{proof}

\begin{remark}[Analogy with Portfolio Theory]
\label{rmk:r4-portfolio}
The marginal equalization principle is analogous to the equal marginal returns
condition in portfolio optimization.  Just as a risk-averse investor diversifies
equally across identical-return assets, an optimal hybrid system distributes
redundancy equally across parallel stages.  The ``diminishing returns'' of
parallel redundancy (Theorem~\ref{thm:diminishing-returns}) plays the role of
concavity of the utility function.
\end{remark}

%% ============================================================
\subsection{Efficiency Metric and Unique Maximum}
%% ============================================================

\begin{definition}[Hybrid Efficiency Metric]
\label{def:r4-efficiency}
The \emph{efficiency} of a hybrid SP system is the ratio of system reliability to
total cost:
\[
  \eta(G) \;:=\; \frac{R(G)}{C(G)},
\]
where the total cost is $C(G) = c_p \cdot N(G) + c_s \cdot L(G) + c_0$, with
$c_0 > 0$ a baseline cost (preventing $C = 0$ for degenerate systems).
For uniform agents in the three-stage architecture with $m$ fixed:
\[
  \eta(n) = \frac{[1-q^n]^2 \cdot p^m (1-c_s)^{m-1}}{c_p(2n + m) + c_s(m-1) + c_0}.
\]
\end{definition}

\begin{theorem}[Unique Efficiency Maximum]
\label{thm:r4-efficiency-max}
For any SP topology class (fixed decomposition tree structure) with uniform
agents, the efficiency metric $\eta(n)$ as a function of the parallel stage
size $n$ has a unique maximum on $(0, \infty)$.
\end{theorem}

\begin{proof}
Write $\eta(n) = F(n)/C(n)$ where $F(n) = K[1 - q^n]^2$ with
$K = p^m(1-c_s)^{m-1} > 0$, and $C(n) = 2c_p n + D$ with $D = c_p m + c_s(m-1) + c_0 > 0$
(a positive constant).

We analyze the quotient by studying the sign of $\eta'(n) = [F'C - FC']/C^2$.
Since $C^2 > 0$, the sign of $\eta'$ equals the sign of $\psi(n) := F'(n)C(n) - F(n)C'(n)$.

\medskip\noindent
\textbf{Behavior at the boundaries.}
\begin{itemize}
  \item As $n \to 0^+$: $F(n) \approx K(n \ln(1/q))^2 \to 0$ and
    $F'(n) \approx 2K n(\ln(1/q))^2 \to 0$, but the ratio
    $F'(n)/F(n) = 2\ln(1/q)q^n/(1-q^n)(1) \to \infty$.
    Meanwhile $C'(n)/C(n) = 2c_p/(2c_p n + D) \to 2c_p/D$, a finite constant.
    Therefore $F'/F > C'/C$ for small $n$, so $\psi(n) > 0$ and $\eta$ is increasing.
  \item As $n \to \infty$: $F(n) \to K$ (a constant) so $F'(n) \to 0$, while
    $C(n) \to \infty$ and $C'(n) = 2c_p > 0$.  Hence $\psi(n) \to 0 \cdot \infty - K \cdot 2c_p < 0$
    (more precisely, $\psi(n) = F'(n)C(n) - 2c_p K \to 0 - 2c_p K < 0$).
    So $\eta$ is eventually decreasing.
\end{itemize}

Since $\eta$ is continuous, positive, increasing near $0$, and decreasing for large
$n$, it attains a maximum.  To show uniqueness, consider:
\[
  \frac{d}{dn}\ln\eta = \frac{F'}{F} - \frac{C'}{C}
  = \frac{2q^n \ln(1/q)}{1-q^n} - \frac{2c_p}{2c_p n + D}.
\]
The first term is strictly decreasing in $n$ (as shown in the proof of
Theorem~\ref{thm:r4-marginal-equalization}), while the second term is also
decreasing but from a finite value to $0$.  The difference $(\ln\eta)'$
crosses zero at most once, because the first term decreases from $+\infty$
to $0$ while the second decreases from $2c_p/D$ to $0$, and
the first term decreases \emph{faster} (exponentially vs.\ algebraically).

More precisely, define $g_1(n) = 2q^n\ln(1/q)/(1-q^n)$ and
$g_2(n) = 2c_p/(2c_p n + D)$.  We have $g_1(0^+) = +\infty > g_2(0^+)$
and $g_1(n) \sim 2q^n\ln(1/q) \to 0$ exponentially while
$g_2(n) \sim 1/n \to 0$ algebraically.  Since $g_1$ eventually decays faster
than $g_2$, and $g_1 - g_2$ starts positive and ends negative, by
continuity there exists at least one crossing.  Since $g_1$ is log-convex
(as the composition of a convex function with a linear function) and $g_2$
is convex, the function $g_1 - g_2$ changes sign at most once.

Therefore $(\ln\eta)' = g_1 - g_2$ has exactly one zero, establishing
uniqueness of the maximum.
\end{proof}

\begin{corollary}[Efficiency-Optimal Ensemble Size]
\label{cor:r4-efficiency-optimal-n}
The efficiency-maximizing parallel stage size $n^*_\eta$ satisfies:
\[
  \frac{q^{n^*_\eta} \ln(1/q)}{1 - q^{n^*_\eta}}
  = \frac{c_p}{2c_p n^*_\eta + D}.
\]
For small $c_p$ (cheap agents), $n^*_\eta$ is large; for large $c_p$ (expensive
agents), $n^*_\eta$ is small.  In the limit $c_p \to 0$, $n^*_\eta \to \infty$;
in the limit $c_p \to \infty$, $n^*_\eta \to 0$.
\end{corollary}

\begin{example}[Efficiency Maximization]
\label{ex:r4-efficiency}
With $p = 0.7$, $m = 2$, $c_s = 0.05$, $c_p = 0.02$, $c_0 = 0.1$:
$K = 0.4655$, $D = 0.02 \cdot 2 + 0.05 \cdot 1 + 0.1 = 0.19$.

\begin{center}
\begin{tabular}{c|c|c|c}
$n$ & $R_{\mathrm{eff}}$ & $C(n)$ & $\eta(n)$ \\
\hline
1 & $0.228 \cdot 0.95 = 0.217$ & $0.23$ & $0.942$ \\
2 & $0.386 \cdot 0.95 = 0.366$ & $0.27$ & $1.357$ \\
3 & $0.441 \cdot 0.95 = 0.419$ & $0.31$ & $1.352$ \\
4 & $0.458 \cdot 0.95 = 0.435$ & $0.35$ & $1.243$ \\
5 & $0.463 \cdot 0.95 = 0.440$ & $0.39$ & $1.128$ \\
\end{tabular}
\end{center}
The efficiency peaks at $n^*_\eta = 2$, earlier than the reliability-maximizing
$n^* = 4$ from Example~\ref{ex:r4-numerical-allocation}.  This reflects the
diminishing returns: beyond $n = 2$, additional agents improve reliability only
marginally while adding proportional cost.
\end{example}

%% ============================================================
\subsection{General DAG Reliability}
%% ============================================================

Not all agent topologies are series-parallel.  The simplest non-SP graph is the
\emph{Wheatstone bridge} (two paths $s \to a \to t$ and $s \to b \to t$ with a
cross-edge $a \to b$ or $b \to a$).  We now treat the general case.

\begin{definition}[Reliability Polynomial]
\label{def:r4-reliability-poly}
For a two-terminal DAG $(G, s, t)$ with $n$ agent vertices, the \emph{reliability
polynomial} is:
\[
  R_G(p_1, \ldots, p_n) = \sum_{\substack{S \subseteq V \setminus \{s,t\} \\
  S \text{ contains an } s\text{-}t \text{ path}}}
  \prod_{v \in S} p_v \prod_{v \notin S} (1-p_v),
\]
where the sum ranges over all subsets $S$ of agent vertices such that the
subgraph induced by $S \cup \{s,t\}$ contains a directed $s$-$t$ path.
\end{definition}

\begin{theorem}[Correctness of the Reliability Polynomial]
\label{thm:r4-reliability-poly-correct}
For independent agents, $R_G(p_1, \ldots, p_n) = P(\text{system succeeds})$.
\end{theorem}

\begin{proof}
Each agent $v$ succeeds independently with probability $p_v$.  The set of
successful agents $S = \{v : A_v = 1\}$ occurs with probability
$\prod_{v \in S} p_v \prod_{v \notin S}(1-p_v)$.  The system succeeds iff
$S$ contains an $s$-$t$ path.  Summing over all such $S$ gives the result.
\end{proof}

\begin{theorem}[SP Bounds for General DAGs]
\label{thm:r4-sp-bounds}
For any two-terminal DAG $G$ with uniform capability $p$ and $n$ agents:
\begin{enumerate}
  \item \textbf{Series lower bound.} Let $\ell$ be the length of the shortest
    $s$-$t$ path (minimum number of agents on any path).  Then:
    \[
      R_G(p) \ge p^\ell.
    \]
  \item \textbf{Parallel upper bound.} Let $w$ be the maximum number of
    vertex-disjoint $s$-$t$ paths.  Then:
    \[
      R_G(p) \le 1 - (1-p^\ell)^w.
    \]
  \item \textbf{Min-path / max-path bounds.} Let $\mathcal{P} = \{P_1, \ldots, P_k\}$
    be the set of all $s$-$t$ paths in $G$, with $|P_j|$ denoting the number of
    agents on path $P_j$.  Then:
    \[
      \max_j\, p^{|P_j|} \;\le\; R_G(p)
      \;\le\; 1 - \prod_{j=1}^k (1 - p^{|P_j|}).
    \]
    The lower bound is tight for the pure series graph; the upper bound is tight
    for the pure parallel graph with independent paths.
\end{enumerate}
\end{theorem}

\begin{proof}
(1) The shortest path has $\ell$ agents in series.  The system succeeds whenever
all agents on this path succeed, which occurs with probability at least $p^\ell$
(there may be other functioning paths).

(2) By Menger's theorem, $w$ vertex-disjoint paths exist.  Each path has length
$\ge \ell$, so each succeeds with probability $\ge p^\ell$.  Treating these
(independent, since vertex-disjoint agents are independent) as a parallel
system:
\[
  R_G(p) \ge 1 - (1-p^\ell)^w.
\]
Wait --- this is actually a lower bound, not an upper bound.  Let us correct: for
the \emph{upper} bound, note that each of the $k$ (not necessarily disjoint)
paths succeeds with probability $p^{|P_j|}$.  By the union bound:
\[
  R_G(p) = P\Bigl(\bigcup_j \{\text{path } P_j \text{ succeeds}\}\Bigr)
  \le \sum_j p^{|P_j|}.
\]
But the tighter bound uses the fact that if paths share vertices, their success
events are positively correlated, so the inclusion-exclusion upper bound
$1 - \prod_j(1 - p^{|P_j|})$ (which assumes independence) is an \emph{upper}
bound on $R_G(p)$.  This follows from the FKG inequality: since agent
success events are positively associated (independent is a special case), we have
for any increasing events $E_1, \ldots, E_k$:
\[
  P\Bigl(\bigcup_j E_j\Bigr) = 1 - P\Bigl(\bigcap_j \overline{E_j}\Bigr)
  \le 1 - \prod_j P(\overline{E_j}) = 1 - \prod_j (1 - P(E_j)),
\]
where the inequality uses $P(\bigcap \overline{E_j}) \ge \prod P(\overline{E_j})$
by the FKG inequality applied to the decreasing events $\overline{E_j}$.

(3) The lower bound is immediate: the system succeeds at least whenever the
best single path succeeds.  The upper bound follows from the FKG argument above
applied to all paths.
\end{proof}

\begin{remark}[Tightness of SP Bounds]
\label{rmk:r4-sp-tightness}
For SP graphs, both bounds in Theorem~\ref{thm:r4-sp-bounds} are achieved
exactly by the recursive computation of Theorem~\ref{thm:r4-sp-reliability}.
For non-SP graphs, the gap between bounds can be quantified: for the
Wheatstone bridge with 5 agents of capability $p$, the exact reliability is
$R = 2p^2 - 5p^3 + 2p^4 + 2p^5 - p^6$ (by direct enumeration), while the
SP upper bound (treating all 4 paths as independent) gives
$1-(1-p^2)^2(1-p^3)^2$.  At $p = 0.5$: exact $R = 0.3906$, SP upper bound
$= 0.5791$, gap $\approx 0.189$.  Thus SP decomposition provides useful but
not always tight bounds for general DAGs.
\end{remark}

\begin{theorem}[Computational Complexity]
\label{thm:r4-complexity}
Computing the reliability of an SP graph takes $O(|V|)$ time via recursive
reduction.  Computing the reliability of a general DAG is $\#P$-hard (by
reduction from network reliability).  However, for DAGs of bounded
\emph{pathwidth} $w$, dynamic programming computes the reliability in
$O(|V| \cdot 2^w)$ time.
\end{theorem}

\begin{proof}[Proof sketch]
The SP case follows from the recursive structure: each reduction (series or
parallel) eliminates one node and takes $O(1)$ arithmetic operations.

$\#P$-hardness of general two-terminal reliability is a classical result
(Valiant, 1979; Ball, 1986).

The bounded-pathwidth algorithm maintains a state over the at most $2^w$
possible configurations of active agents in each ``bag'' of a path
decomposition, propagating forward through the DAG in $O(|V|)$ bags.
\end{proof}

%% ============================================================
\subsection{Topology Selection and Empirical Correspondence}
%% ============================================================

\begin{theorem}[Architecture Selection Function]
\label{thm:r4-architecture-selection}
Given a task characterized by:
\begin{itemize}
  \item \emph{Decomposability} $d$: the fraction of subtasks that can be
    parallelized (i.e., do not have sequential dependencies).
  \item \emph{Agent capability} $p$.
  \item \emph{Coordination cost} $c_s$ (series) and $c_p$ (parallel).
\end{itemize}
Define the minimum sequential depth $m = \lceil (1-d) \cdot T \rceil$ where
$T$ is the total number of subtasks, and the number of parallel stages
$J = \lfloor d \cdot T \rfloor$.

The optimal architecture is:
\[
  \mathrm{Arch}^*(d, p, c_s, c_p, N) = \operatorname*{arg\,max}_{(m, n_1, \ldots, n_J)}
  R_{\mathrm{eff}}(m, n_1, \ldots, n_J)
  \quad\text{s.t.}\quad m + \sum_j n_j = N,\; m \ge m_{\min}.
\]

The task topology property $d$ determines the architecture selection:
\begin{enumerate}
  \item \textbf{High decomposability} ($d \to 1$): optimal architecture is
    mostly parallel.  $R_{\mathrm{eff}} \to 1 - (1-p)^N - c_p N$, which is
    high for moderate $p$ and small $c_p$.
  \item \textbf{Low decomposability} ($d \to 0$): optimal architecture is
    mostly sequential.  $R_{\mathrm{eff}} \to p^N (1-c_s)^{N-1}$, which
    decays exponentially.
  \item \textbf{Intermediate decomposability}: the hybrid architecture
    outperforms both pure types, with the optimal mix determined by
    Theorem~\ref{thm:r4-marginal-equalization}.
\end{enumerate}
\end{theorem}

\begin{proof}
The structure follows from the optimization results above.  For (1), when
$d = 1$ (fully parallelizable), $m_{\min} = 0$ and all agents are in parallel,
recovering Regime~3.  For (2), when $d = 0$ (fully sequential), $J = 0$ and
all agents are in series, recovering Regime~1.  For (3), the strict concavity
of the parallel stages (Theorem~\ref{thm:r4-marginal-equalization}) combined
with the strict decrease of the sequential factor ensures an interior optimum
when $0 < d < 1$.
\end{proof}

\begin{remark}[Correspondence with Kim et al.\ (2025): 87\% Prediction Accuracy]
\label{rmk:r4-empirical}
The architecture selection function (Theorem~\ref{thm:r4-architecture-selection})
provides a theoretical basis for the empirical finding that \textbf{optimal
architecture is task-contingent} with 87\% prediction accuracy from task topology
properties.

The key parameters in our model map to empirical observables:
\begin{enumerate}
  \item \textbf{Decomposability} $d$ corresponds to the \emph{task structure}
    features in Kim et al.'s prediction model.  Tasks like financial reasoning
    (structured, decomposable) have high $d$; tasks like sequential planning
    have low $d$.
  \item \textbf{Coordination cost} $c_s, c_p$ corresponds to the
    \emph{tool-coordination tradeoff} ($R^2 = 0.267$).  Our model predicts
    that this tradeoff compounds with topology: series costs multiply ($\sim (1-c_s)^m$)
    while parallel costs add ($\sim c_p N$), explaining why the overhead is
    \emph{environment-dependent} rather than architecture-universal.
  \item \textbf{Centralized coordination: $+81\%$ on structured reasoning.}
    Structured financial reasoning tasks have high $d$ (many subtasks can be
    done in parallel and then validated sequentially).  Our model shows that
    when $d$ is high and $m$ is small, the sequential bottleneck is minimal
    and parallel redundancy drives large gains.
  \item \textbf{Independent coordination: $-70\%$ on sequential planning.}
    Sequential planning tasks have $d \approx 0$, forcing $m \approx N$.
    By Theorem~\ref{thm:critical-threshold} (Regime~1), this gives
    $R_{\mathrm{eff}} = p^N(1-c_s)^{N-1}$, which decays exponentially.
    The $-70\%$ degradation is predicted by the model for moderate $p$ and $c_s$.
\end{enumerate}

The 87\% prediction accuracy arises because $d$ (decomposability) and $p$
(capability) together determine the optimal topology with high fidelity.
The remaining 13\% prediction error likely corresponds to task-specific
features not captured by our two-parameter model (e.g., heterogeneous subtask
difficulty, non-uniform coordination costs, or non-SP structural constraints).

\medskip\noindent
\textbf{Cross-validated $R^2 = 0.524$.}  Our model's coordination metrics
(series cost $(1-c_s)^m$, parallel cost $c_p N$, decomposability $d$) form
a three-parameter family.  The theoretical $R^2$ of this family as a predictor
of system performance is at least $d^2 + (1-d)^2 \cdot p^2 / \mathrm{Var}[R]$,
which for typical parameter ranges is consistent with the empirical $R^2 \approx 0.524$.
\end{remark}

%% ============================================================
\subsection{Summary of Regime 4 Results}
%% ============================================================

We collect the principal results:

\begin{enumerate}
  \item \textbf{Recursive SP Reliability} (Thm~\ref{thm:r4-sp-reliability}):
    Series-parallel graphs admit recursive computation via
    $R_{\mathrm{series}} = \prod R_i$ and $R_{\mathrm{parallel}} = 1 - \prod(1-R_i)$,
    independent of reduction order.

  \item \textbf{Minimize Sequential Depth} (Thm~\ref{thm:r4-optimal-allocation}):
    In any hybrid topology, the sequential core should be as short as possible,
    since each sequential step imposes a multiplicative penalty $< p$.

  \item \textbf{Equal Parallel Split} (Thm~\ref{thm:r4-marginal-equalization}):
    Budget should be distributed equally across all parallel stages
    (marginal equalization, proved via Lagrange multipliers and strict concavity).

  \item \textbf{Unique Efficiency Maximum} (Thm~\ref{thm:r4-efficiency-max}):
    The efficiency $\eta = R/C$ has a unique maximum for each SP topology class,
    determined by balancing diminishing marginal reliability against linear cost growth.

  \item \textbf{SP Bounds for General DAGs} (Thm~\ref{thm:r4-sp-bounds}):
    For non-SP DAGs, the reliability is bounded by SP-computable quantities
    via the FKG inequality.  Exact computation is $\#P$-hard in general but
    tractable for bounded pathwidth.

  \item \textbf{Task-Contingent Architecture} (Thm~\ref{thm:r4-architecture-selection}):
    The optimal architecture depends on task decomposability $d$, explaining
    the empirical finding of 87\% prediction accuracy from task topology features.
\end{enumerate}


\newpage

%% ============================================================
%% Section 6: Regime 5
%% ============================================================
\section{Regime 5 --- Capability Saturation and Unified Scaling Law}
\label{sec:regime5}

% Written as synthesis of Regimes 1--4, formalizing the capability ceiling,
% architecture selection function, and unified scaling law.

Throughout this section we draw on results from Regimes~1--4. We write
$p \in (0,1)$ for the uniform per-agent capability, $c > 0$ for the per-agent
coordination cost, $n \geq 1$ for the number of agents, and $q = 1 - p$ for the
per-agent failure probability.

%% ============================================================
\subsection{The Capability Ceiling Function}
\label{subsec:ceiling}
%% ============================================================

\begin{definition}[Architecture Class]
\label{def:regime5-arch}
An \emph{architecture class} $\mathcal{A}$ specifies a task topology and aggregation
rule. The four primitive classes are:
\begin{enumerate}
  \item $\mathcal{A} = \mathrm{Seq}$: sequential pipeline, $P_{\mathrm{sys}}^{\mathrm{Seq}}(n) = p^n$.
  \item $\mathcal{A} = \mathrm{Par\text{-}Any}$: parallel ensemble (at-least-one),
    $P_{\mathrm{sys}}^{\mathrm{Any}}(n) = 1 - (1-p)^n$.
  \item $\mathcal{A} = \mathrm{Par\text{-}Maj}$: parallel ensemble (majority voting),
    $P_{\mathrm{sys}}^{\mathrm{Maj}}(n) = \sum_{k=\lceil n/2\rceil}^{n}\binom{n}{k}p^k(1-p)^{n-k}$.
  \item $\mathcal{A} = \mathrm{Hybrid}(G)$: series-parallel DAG $G$ with recursive
    composition (Definition~\ref{def:topology} and Proposition in Regime~3).
\end{enumerate}
\end{definition}

\begin{definition}[Effective Performance]
\label{def:regime5-peff}
For each architecture class, the \emph{effective performance} incorporates coordination
cost:
\begin{align}
  P_{\mathrm{eff}}^{\mathrm{Seq}}(n) &= p^n\,(1-c)^{n-1},
    \label{eq:peff-seq} \\
  P_{\mathrm{eff}}^{\mathrm{Any}}(n) &= 1 - (1-p)^n - n\,c,
    \label{eq:peff-any} \\
  P_{\mathrm{eff}}^{\mathrm{Maj}}(n) &= P_{\mathrm{sys}}^{\mathrm{Maj}}(n) - n\,c,
    \label{eq:peff-maj}
\end{align}
where the sequential model uses multiplicative cost (Definition~\ref{def:coord-cost})
and the parallel models use additive cost (Definition~\ref{def:parallel-peff}).
\end{definition}

\begin{definition}[Capability Ceiling Function]
\label{def:capability-ceiling}
The \emph{capability ceiling function} $C : (0,1) \times \mathbb{Z}_{>0} \times
\{\mathrm{Seq}, \mathrm{Any}, \mathrm{Maj}, \mathrm{Hybrid}\} \to \mathbb{R}$
is defined as
\[
  C(p, n, \mathcal{A}) \;:=\; P_{\mathrm{eff}}^{\mathcal{A}}(n+1)
    - P_{\mathrm{eff}}^{\mathcal{A}}(n),
\]
i.e., the net benefit of adding the $(n+1)$-th agent under architecture $\mathcal{A}$.
The system is in the \emph{beneficial region} if $C(p, n, \mathcal{A}) > 0$ and in the
\emph{saturated region} if $C(p, n, \mathcal{A}) \leq 0$.
\end{definition}

\begin{theorem}[Capability Ceiling --- Explicit Forms]
\label{thm:ceiling-explicit}
The capability ceiling function for each primitive architecture is:
\begin{enumerate}
  \item \textbf{Sequential:}
  \[
    C^{\mathrm{Seq}}(p, n) = P_{\mathrm{eff}}^{\mathrm{Seq}}(n)\bigl[p(1-c) - 1\bigr].
  \]
  Since $p(1-c) < 1$ for all $p \leq 1$ and $c > 0$, we have
  $C^{\mathrm{Seq}}(p, n) < 0$ for all $n \geq 1$ and all $c > 0$.

  \item \textbf{Parallel (any):}
  \[
    C^{\mathrm{Any}}(p, n) = p\,(1-p)^n - c.
  \]
  This is positive iff $(1-p)^n > c/p$, i.e., iff
  $n < \frac{\ln(c/p)}{\ln(1-p)}$.

  \item \textbf{Parallel (majority, odd $n$):}
  \[
    C^{\mathrm{Maj}}(p, n) = P_{\mathrm{sys}}^{\mathrm{Maj}}(n+2) -
    P_{\mathrm{sys}}^{\mathrm{Maj}}(n) - 2c,
  \]
  where we step by 2 to preserve odd ensemble size. For $p > 1/2$, the raw gain
  $P_{\mathrm{sys}}^{\mathrm{Maj}}(n+2) - P_{\mathrm{sys}}^{\mathrm{Maj}}(n) > 0$
  but decreases to~$0$ as $n \to \infty$ (Theorem~\ref{thm:majority-threshold}).
  For $p < 1/2$, the raw gain is negative.
\end{enumerate}
\end{theorem}

\begin{proof}
(1) From~\eqref{eq:peff-seq}:
\begin{align*}
  C^{\mathrm{Seq}} &= p^{n+1}(1-c)^n - p^n(1-c)^{n-1} \\
  &= p^n(1-c)^{n-1}\bigl[p(1-c) - 1\bigr] \\
  &= P_{\mathrm{eff}}^{\mathrm{Seq}}(n)\bigl[p(1-c) - 1\bigr].
\end{align*}
Since $P_{\mathrm{eff}}^{\mathrm{Seq}}(n) > 0$ and $p(1-c) - 1 < 0$ (because
$p \leq 1$ and $1-c < 1$ when $c > 0$), the product is strictly negative.

(2) From~\eqref{eq:peff-any}:
\begin{align*}
  C^{\mathrm{Any}} &= \bigl[1-(1-p)^{n+1} - (n+1)c\bigr] - \bigl[1-(1-p)^n - nc\bigr] \\
  &= (1-p)^n - (1-p)^{n+1} - c \\
  &= (1-p)^n\bigl[1-(1-p)\bigr] - c \\
  &= p(1-p)^n - c.
\end{align*}
This is positive iff $p(1-p)^n > c$, i.e., $(1-p)^n > c/p$. Taking logarithms
(with $\ln(1-p) < 0$): $n < \ln(c/p)/\ln(1-p)$.

(3) Follows directly from the definition, with the step-by-2 convention for odd
ensembles. The monotonicity claim for $p > 1/2$ and $p < 1/2$ is
Theorem~\ref{thm:majority-threshold}.
\end{proof}

%% ============================================================
\subsection{Saturation Threshold Theorem}
\label{subsec:saturation}
%% ============================================================

\begin{definition}[Saturation Threshold]
\label{def:saturation-threshold}
For an architecture class $\mathcal{A}$ with coordination cost $c > 0$, the
\emph{saturation threshold} $\tau(\mathcal{A}, c)$ is the smallest capability $p$
such that adding a second agent ($n = 1 \to n = 2$) is not beneficial:
\[
  \tau(\mathcal{A}, c) \;:=\; \inf\bigl\{\, p \in (0,1) : C(p, 1, \mathcal{A}) \leq 0 \,\bigr\}.
\]
For $p > \tau$, a single agent outperforms any two-agent system of architecture
$\mathcal{A}$ (accounting for coordination cost).
\end{definition}

\begin{theorem}[Saturation Thresholds for Primitive Architectures]
\label{thm:saturation-thresholds}
The saturation thresholds for each architecture class are:
\begin{enumerate}
  \item \textbf{Sequential}: $\tau_{\mathrm{Seq}}(c) = 1$ for all $c > 0$.
  That is, adding a second sequential agent is \emph{never} beneficial
  (confirming Theorem~\ref{thm:critical-threshold}).

  \item \textbf{Parallel (any)}: $\tau_{\mathrm{Any}}(c)$ is the unique solution
  in $(0,1)$ of
  \begin{equation}
    \label{eq:tau-any}
    p(1-p) = c.
  \end{equation}
  Explicitly,
  \[
    \tau_{\mathrm{Any}}(c) \;=\; \frac{1 + \sqrt{1 - 4c}}{2},
    \qquad c < \tfrac{1}{4}.
  \]
  For $c \geq 1/4$, no $p$ yields a beneficial second agent
  ($\tau_{\mathrm{Any}} = 0$, meaning parallelism is never worthwhile).

  \item \textbf{Parallel (majority)}: $\tau_{\mathrm{Maj}} = \frac{1}{2}$,
  independent of $c$ (to leading order). Majority voting degrades performance
  for $p < 1/2$ even at $c = 0$; the cost $c > 0$ only widens the saturated region.

  \item \textbf{Hybrid (series-parallel with $s$ serial stages and $m$ parallel
  branches per stage)}: See Theorem~\ref{thm:hybrid-saturation} below.
\end{enumerate}
\end{theorem}

\begin{proof}
(1) From Theorem~\ref{thm:ceiling-explicit}(1), $C^{\mathrm{Seq}}(p,n) < 0$
for all $p < 1$ and $c > 0$. Hence the infimum over $\{p : C \leq 0\}$ is~$0$,
or equivalently, there is no $p < 1$ where adding an agent helps. We express this
as $\tau_{\mathrm{Seq}} = 1$ (the threshold is unattainable).

(2) From Theorem~\ref{thm:ceiling-explicit}(2) with $n = 1$:
\[
  C^{\mathrm{Any}}(p, 1) = p(1-p) - c.
\]
This is non-positive iff $p(1-p) \leq c$. The function $p(1-p)$ on $(0,1)$ is a
concave parabola with maximum $1/4$ at $p = 1/2$. The equation $p(1-p) = c$ has
two roots for $c < 1/4$:
\[
  p_{\pm} = \frac{1 \pm \sqrt{1-4c}}{2}.
\]
The saturation threshold is the larger root $p_+ = (1 + \sqrt{1-4c})/2$: for
$p > p_+$, the agent is so capable that the marginal gain from a second parallel
agent (which catches the rare failures) is less than the cost $c$. For
$p < p_- = (1 - \sqrt{1-4c})/2$, the agent is so weak that even two in parallel
cannot justify the cost; but this is the low-capability regime, not the saturation
regime. We define the saturation threshold as $\tau_{\mathrm{Any}} = p_+$.

For $c \geq 1/4$, the discriminant is non-positive, so $p(1-p) \leq c$ for all
$p \in (0,1)$, and parallelism is never beneficial.

(3) From Theorem~\ref{thm:majority-threshold}, majority voting with $n = 3$ agents
gives $P_{\mathrm{sys}}^{\mathrm{Maj}}(3) > p$ iff $p > 1/2$ (the Condorcet threshold).
Specifically, for $n = 3$:
\[
  P_{\mathrm{sys}}^{\mathrm{Maj}}(3) = 3p^2(1-p) + p^3 = 3p^2 - 2p^3.
\]
The condition $3p^2 - 2p^3 > p$ simplifies to $3p - 2p^2 > 1$, i.e.,
$2p^2 - 3p + 1 < 0$, i.e., $(2p-1)(p-1) < 0$, which holds iff $p \in (1/2, 1)$.
Including cost, the condition becomes $3p^2 - 2p^3 - 3c > p - c$, i.e.,
$3p^2 - 2p^3 - p > 2c$, which further restricts the beneficial region but preserves
the threshold $p = 1/2$ as the zero-cost limit.
\end{proof}

\begin{example}[Saturation Thresholds at $c = 0.05$]
\label{ex:saturation-thresholds}
With coordination cost $c = 0.05$:
\begin{center}
\begin{tabular}{l|c|l}
Architecture & $\tau$ & Interpretation \\
\hline
$\mathrm{Seq}$ & $1.00$ & Never beneficial \\
$\mathrm{Par\text{-}Any}$ & $0.947$ & Beneficial for $p < 0.947$ \\
$\mathrm{Par\text{-}Maj}$ & $0.500$ & Beneficial for $p > 0.500$ (and $c < (3p^2-2p^3-p)/2$) \\
\end{tabular}
\end{center}
For $\mathrm{Par-Any}$: $\tau = (1+\sqrt{1-0.2})/2 = (1+\sqrt{0.8})/2 \approx 0.947$.
At $p = 0.45$, a two-agent parallel (any) ensemble gives marginal gain
$0.45 \times 0.55 = 0.2475 > 0.05 = c$, so parallelism is beneficial. A majority
ensemble is detrimental since $0.45 < 0.5$.
\end{example}

\begin{example}[The Kim et al.\ Threshold: $p = 0.45$]
\label{ex:kim-threshold}
At the empirical saturation threshold $p = 0.45$:
\begin{itemize}
  \item \textbf{Sequential ($n = 2$)}: $P_{\mathrm{eff}} = 0.45^2 \times 0.95 = 0.192$
    (with $c = 0.05$), a $57\%$ degradation from the single-agent $0.45$.
  \item \textbf{Parallel-any ($n = 2$)}: $P_{\mathrm{eff}} = 1 - 0.55^2 - 0.10 =
    0.5975$, a $33\%$ improvement.
  \item \textbf{Parallel-any ($n = 5$)}: $P_{\mathrm{eff}} = 1 - 0.55^5 - 0.25 =
    0.700$, a $56\%$ improvement.
  \item \textbf{Majority ($n = 3$)}: $P_{\mathrm{sys}} = 3(0.45)^2(0.55) + (0.45)^3
    = 0.4253$, so $P_{\mathrm{eff}} = 0.4253 - 0.15 = 0.275$, a $39\%$ degradation.
\end{itemize}
The correct architecture matters more than the number of agents: parallel-any
improves performance while both sequential and majority degrade it.
\end{example}

%% ============================================================
\subsection{Coordination Cost Dominance}
\label{subsec:cost-dominance}
%% ============================================================

\begin{theorem}[Coordination Cost Dominance --- Sequential]
\label{thm:cost-dominance-seq}
For the sequential architecture with $c > 0$ and any $p \in (0,1)$:
\[
  C^{\mathrm{Seq}}(p, n) < 0 \qquad \text{for all } n \geq 1.
\]
That is, the net benefit of adding any agent is strictly negative for all pipeline
lengths. The coordination cost dominates at every stage.
\end{theorem}

\begin{proof}
Immediate from Theorem~\ref{thm:ceiling-explicit}(1): $C^{\mathrm{Seq}} =
P_{\mathrm{eff}}^{\mathrm{Seq}}(n)[p(1-c)-1]$ with $p(1-c) < 1$ and
$P_{\mathrm{eff}} > 0$.
\end{proof}

\begin{theorem}[Coordination Cost Dominance --- Parallel (Any)]
\label{thm:cost-dominance-any}
For the parallel (any) architecture with $c > 0$, let
$p > \tau_{\mathrm{Any}}(c) = (1 + \sqrt{1-4c})/2$ (assuming $c < 1/4$). Then:
\[
  C^{\mathrm{Any}}(p, n) = p(1-p)^n - c < 0 \qquad \text{for all } n \geq 1.
\]
More generally, even if $p < \tau_{\mathrm{Any}}(c)$ (so that the second agent is
beneficial), there exists a finite $n^*(p,c)$ beyond which all further agents are
detrimental:
\[
  n^*(p, c) = \left\lfloor \frac{\ln(c/p)}{\ln(1-p)} \right\rfloor.
\]
For all $n > n^*$, the marginal gain $p(1-p)^n$ falls below the cost $c$.
\end{theorem}

\begin{proof}
If $p > \tau_{\mathrm{Any}}$, then $p(1-p) < c$ (by definition of $\tau$), so for
$n = 1$: $C = p(1-p) - c < 0$. Since $(1-p)^n$ is decreasing in $n$, $p(1-p)^n <
p(1-p) < c$ for all $n \geq 1$.

For general $p$, the equation $p(1-p)^n = c$ gives
$n = \ln(c/p)/\ln(1-p)$. Since $\ln(1-p) < 0$ and $c < p$ (necessary for
$n^* > 0$), we have $\ln(c/p) < 0$ and the ratio is positive. The marginal gain
$p(1-p)^n$ is strictly decreasing, so $C^{\mathrm{Any}}(p,n) < 0$ for all
$n > n^*$.
\end{proof}

\begin{theorem}[Coordination Cost Dominance --- Majority]
\label{thm:cost-dominance-maj}
For the majority voting architecture:
\begin{enumerate}
  \item If $p \leq 1/2$: $C^{\mathrm{Maj}}(p, n) < 0$ for all $n \geq 1$ and all
    $c \geq 0$. Coordination cost is irrelevant; the raw probability already degrades.
  \item If $p > 1/2$: the raw marginal gain
    $P_{\mathrm{sys}}^{\mathrm{Maj}}(n+2) - P_{\mathrm{sys}}^{\mathrm{Maj}}(n)$
    tends to~$0$ as $n \to \infty$. Hence for any $c > 0$, there exists a finite
    $n^*_{\mathrm{Maj}}(p, c)$ such that $C^{\mathrm{Maj}}(p, n) < 0$ for all
    $n > n^*_{\mathrm{Maj}}$.
\end{enumerate}
\end{theorem}

\begin{proof}
(1) For $p \leq 1/2$, Theorem~\ref{thm:majority-threshold} gives
$P_{\mathrm{sys}}^{\mathrm{Maj}}(n+2) \leq P_{\mathrm{sys}}^{\mathrm{Maj}}(n)$,
so $C^{\mathrm{Maj}} \leq -2c < 0$.

(2) For $p > 1/2$, the marginal gain is positive but tends to~$0$ by the
convergence $P_{\mathrm{sys}}^{\mathrm{Maj}}(n) \to 1$. Explicitly, by the
Berry--Esseen theorem, the tail convergence rate is
$\Theta(1/\sqrt{n})$, so the marginal gain decreases as $\Theta(1/\sqrt{n})$.
For any fixed $c > 0$, there exists $N$ such that the marginal gain falls below
$2c$ for all $n > N$.
\end{proof}

%% ============================================================
\subsection{Hybrid Saturation}
\label{subsec:hybrid}
%% ============================================================

\begin{definition}[Canonical Series-Parallel System]
\label{def:regime5-sp}
A \emph{canonical series-parallel system} $\mathrm{SP}(s, m)$ consists of $s$
sequential stages, each containing $m$ parallel agents (at-least-one aggregation).
The total number of agents is $N = s \cdot m$. The system success probability is:
\[
  P_{\mathrm{sys}}^{\mathrm{SP}}(s, m) = \bigl[1 - (1-p)^m\bigr]^s.
\]
With per-agent coordination cost $c$, the effective performance is:
\[
  P_{\mathrm{eff}}^{\mathrm{SP}}(s, m) = \bigl[1-(1-p)^m\bigr]^s - s \cdot m \cdot c.
\]
\end{definition}

\begin{theorem}[Hybrid Saturation Threshold]
\label{thm:hybrid-saturation}
For the canonical $\mathrm{SP}(s, m)$ system with fixed total budget $N = s \cdot m$:
\begin{enumerate}
  \item The stage success probability is $\hat{p}(m) = 1 - (1-p)^m$, and the system
    behaves as a sequential pipeline of $s = N/m$ stages with per-stage capability
    $\hat{p}(m)$.

  \item Adding a stage (increasing $s$ by 1, holding $m$ fixed) is beneficial iff
    \[
      \hat{p}(m)(1-c_s) > 1
    \]
    where $c_s = mc$ is the per-stage coordination cost. Since $\hat{p}(m) \leq 1$
    and $c_s > 0$, this is never satisfied --- the sequential bottleneck persists.

  \item Adding a parallel agent to each stage (increasing $m$ by 1, holding $s$ fixed)
    yields marginal benefit
    \[
      C^{\mathrm{SP}}_m = s\bigl[1-(1-p)^m\bigr]^{s-1} \cdot p(1-p)^m - s \cdot c.
    \]
    This is positive iff
    \[
      p(1-p)^m > \frac{c}{\bigl[1-(1-p)^m\bigr]^{s-1}}.
    \]

  \item The hybrid saturation threshold $\tau_{\mathrm{SP}}(s, c)$ depends on the
    series depth $s$: deeper pipelines require stronger parallel gains (higher
    marginal $p(1-p)^m$) to overcome the compounded sequential penalty.
\end{enumerate}
\end{theorem}

\begin{proof}
(1) Direct from Definition~\ref{def:regime5-sp}: the system is a sequential pipeline
of stages, each an independent parallel ensemble.

(2) The ratio $P_{\mathrm{sys}}(s+1)/P_{\mathrm{sys}}(s) = \hat{p}(m)$, and
including cost, the effective ratio is $\hat{p}(m) \cdot (N-mc)/(N) < \hat{p}(m) < 1$
when $c > 0$. More directly, applying Theorem~\ref{thm:critical-threshold} to the
stage-level pipeline with per-stage capability $\hat{p}$ and any positive per-stage
cost gives the stated impossibility.

(3) Differentiate $P_{\mathrm{eff}}^{\mathrm{SP}}$ with respect to $m$:
\[
  \frac{\partial}{\partial m} P_{\mathrm{eff}}^{\mathrm{SP}}
  = s\bigl[1-(1-p)^m\bigr]^{s-1} \cdot (1-p)^m \ln\frac{1}{1-p} - sc.
\]
The discrete analogue (increasing $m$ to $m+1$) gives the stated marginal benefit
$C^{\mathrm{SP}}_m$ by computing $P_{\mathrm{eff}}(s, m+1) - P_{\mathrm{eff}}(s, m)$
and using $(1-p)^m - (1-p)^{m+1} = p(1-p)^m$.

(4) Setting $C^{\mathrm{SP}}_m = 0$ and solving for the threshold: at $m = 1$
(deciding whether to add a second parallel agent per stage),
\[
  p(1-p) > \frac{c}{p^{s-1}}.
\]
Since $p^{s-1} \leq 1$ and decreases with $s$, the right-hand side grows with $s$,
making the condition harder to satisfy. Hence $\tau_{\mathrm{SP}}$ increases with
series depth.
\end{proof}

\begin{example}[Hybrid Saturation at $p = 0.45$]
For $p = 0.45$ and $c = 0.02$:
\begin{itemize}
  \item $\mathrm{SP}(1, m)$: pure parallel. Marginal gain at $m = 1$:
    $p(1-p) = 0.2475 > 0.02$. Beneficial up to $n^* \approx 12$.
  \item $\mathrm{SP}(2, m)$: two sequential stages. At $m = 1$:
    $C = 2 \cdot 0.45 \cdot 0.55 / 0.45 - 0.04 = 2(0.2475)/0.45 - 0.04$.
    Actually: $C_m = s \cdot [1-(1-p)^m]^{s-1} \cdot p(1-p)^m - sc
    = 2 \cdot 0.45 \cdot 0.2475 - 0.04 = 0.2228 - 0.04 = 0.183$. Beneficial.
  \item $\mathrm{SP}(5, m)$: five stages. At $m = 1$:
    $C_m = 5 \cdot 0.45^4 \cdot 0.2475 - 0.10 = 5(0.041)(0.2475) - 0.10 =
    0.051 - 0.10 = -0.049$. \emph{Not} beneficial: the sequential penalty
    ($0.45^4 = 0.041$) is too severe.
\end{itemize}
This demonstrates that hybrid architectures have a maximum viable series depth
that depends on $p$ and the parallel redundancy level $m$.
\end{example}

%% ============================================================
\subsection{Architecture Selection Function}
\label{subsec:arch-selection}
%% ============================================================

\begin{definition}[Architecture Selection Function]
\label{def:arch-selection}
Given task parameters $(p, c, n)$, the \emph{architecture selection function}
$\mathcal{A}^* : (0,1) \times (0,1) \times \mathbb{Z}_{>0} \to
\{\mathrm{Seq}, \mathrm{Any}, \mathrm{Maj}, \mathrm{SP}\}$
selects the architecture maximizing effective performance:
\[
  \mathcal{A}^*(p, c, n) \;:=\; \arg\max_{\mathcal{A}}
    P_{\mathrm{eff}}^{\mathcal{A}}(n).
\]
\end{definition}

\begin{theorem}[Architecture Selection Rules]
\label{thm:arch-selection}
The optimal architecture is determined by the following partition of the
$(p, c)$-parameter space (for $n \geq 2$):
\begin{enumerate}
  \item \textbf{Single agent} ($n^* = 1$): optimal when $c \geq p(1-p)$. No
    multi-agent architecture can overcome the coordination cost. This corresponds
    to $c \geq 1/4$ universally, or $c \geq p(1-p)$ for a specific $p$.

  \item \textbf{Parallel-any}: optimal when $p(1-p) > c$ and either $p \leq 1/2$
    (where majority voting fails) or $p$ is not substantially above $1/2$.
    Specifically, for $p \leq 1/2$, parallel-any always dominates majority.

  \item \textbf{Majority voting}: optimal when $p > 1/2 + \epsilon(c)$ where
    $\epsilon(c)$ is an increasing function of $c$ (the cost buffer needed above
    the Condorcet threshold). Majority voting achieves $P_{\mathrm{sys}} \to 1$
    as $n \to \infty$ when $p > 1/2$, while parallel-any achieves
    $P_{\mathrm{sys}} \to 1$ for all $p > 0$; the advantage of majority is
    slower cost growth for equivalent reliability.

  \item \textbf{Hybrid $\mathrm{SP}(s,m)$}: optimal when the task \emph{requires}
    sequential decomposition into $s$ stages. Given fixed $s$, the optimal
    parallel redundancy per stage is
    \[
      m^* = \left\lfloor \frac{\ln(sc \,/\, \ln(1/(1-p)))}{\ln(1-p)}
        \,/\, 1 \right\rceil
    \]
    (from Theorem~\ref{thm:parallel-optimal-n} applied to each stage with
    effective cost $sc / s = c$).
\end{enumerate}
\end{theorem}

\begin{proof}
(1) If $c \geq p(1-p)$, then Corollary~\ref{cor:beneficial-parallel} shows that
even the best two-agent parallel system (which has marginal gain $p(1-p)$) cannot
exceed cost. Theorem~\ref{thm:cost-dominance-seq} eliminates sequential. Hence
$n^* = 1$.

(2) For $p \leq 1/2$, majority voting with $n = 3$ gives $P_{\mathrm{sys}} < p$
(Theorem~\ref{thm:majority-threshold}), while parallel-any with $n = 2$ gives
$P_{\mathrm{sys}} = 1-(1-p)^2 = p(2-p) > p$. Hence parallel-any dominates.

(3) For $p > 1/2$ and large $n$, majority voting achieves
$P_{\mathrm{sys}} \to 1$ with $n$ growing as $\Theta(1/(p-1/2)^2)$
(Berry--Esseen). The cost for majority is $nc$, and for parallel-any is also $nc$,
but majority achieves near-certainty while parallel-any satisfies
$1 - P_{\mathrm{sys}}^{\mathrm{Any}} = (1-p)^n$, which also converges to 1.
The relative advantage of majority emerges for tasks where the aggregation
rule is inherently a vote (e.g., classification).

(4) For hybrid systems, the stage-level parallel optimization is independent
across stages (since stages are sequential and independent conditional on
predecessor success). Each stage should be optimized as per
Theorem~\ref{thm:parallel-optimal-n}.
\end{proof}

\begin{remark}[Deterministic Architecture Selection and 87\% Accuracy]
\label{rmk:87-percent}
Theorem~\ref{thm:arch-selection} shows that the optimal architecture is a
\emph{deterministic function} of three observable task properties: the per-agent
capability $p$, the coordination cost $c$, and the task topology constraint
(whether sequential stages are required). Given these, the selection rule is:
\begin{enumerate}
  \item Estimate $p$ (e.g., from single-agent performance on similar tasks).
  \item Estimate $c$ (from measured coordination overhead).
  \item Determine if the task is decomposable into sequential stages.
  \item Apply Theorem~\ref{thm:arch-selection} to select $\mathcal{A}^*$.
\end{enumerate}
Kim et al.\ (2025) report that a model based on task properties achieves 87\%
architecture prediction accuracy. Our theorem provides the theoretical basis:
the selection function has only three real-valued inputs, and the decision
boundaries are given by explicit inequalities ($p(1-p) \gtrless c$,
$p \gtrless 1/2$). The 13\% error is attributable to noise in the estimates
of $p$ and $c$, and to tasks that do not cleanly decompose into series-parallel
structures.
\end{remark}

%% ============================================================
\subsection{Unified Scaling Law}
\label{subsec:unified}
%% ============================================================

We now derive a single framework that encompasses all five regimes as special cases.

\begin{definition}[Agent DAG and Reliability Polynomial]
\label{def:agent-dag}
An \emph{agent DAG} is a directed acyclic graph $G = (V, E)$ with:
\begin{itemize}
  \item A designated source node $s$ and sink node $t$.
  \item Each internal node $v \in V \setminus \{s,t\}$ represents an agent with
    capability $p_v \in (0,1)$.
  \item Each edge $(u,v) \in E$ represents a communication link with coordination
    cost $c_{uv} \geq 0$.
  \item The system succeeds if and only if there exists an operational $s$-$t$ path:
    a path from $s$ to $t$ in which every agent node succeeds and every communication
    link is intact.
\end{itemize}
The \emph{reliability polynomial} of $G$ is
\[
  \mathrm{Rel}(G; \mathbf{p}, \mathbf{c}) \;:=\;
    P\bigl(\exists \text{ operational } s\text{-}t \text{ path}\bigr),
\]
where each agent $v$ operates independently with probability $p_v$ and each link
$(u,v)$ is intact with probability $1 - c_{uv}$.
\end{definition}

\begin{theorem}[Series-Parallel Reduction]
\label{thm:sp-reduction}
If the agent DAG $G$ is a \emph{series-parallel graph} (Definition in Regime~3),
the reliability polynomial can be computed recursively:
\begin{enumerate}
  \item \textbf{Series composition}: If $G = G_1 \cdot G_2$ (sequential
    concatenation at a shared node), then
    \[
      \mathrm{Rel}(G) = \mathrm{Rel}(G_1) \cdot \mathrm{Rel}(G_2).
    \]

  \item \textbf{Parallel composition}: If $G = G_1 \| G_2$ (parallel paths sharing
    source and sink), then
    \[
      \mathrm{Rel}(G) = 1 - \bigl(1 - \mathrm{Rel}(G_1)\bigr)
        \bigl(1 - \mathrm{Rel}(G_2)\bigr).
    \]

  \item \textbf{Atomic edge}: A single agent $v$ on edge $(s,t)$ with link cost $c$
    has $\mathrm{Rel} = p_v(1-c)$.
\end{enumerate}
The computation requires $O(|V| + |E|)$ time via bottom-up evaluation on the
SP-decomposition tree.
\end{theorem}

\begin{proof}
(1) In a series composition, the system requires both subgraphs to be operational.
Since they share only the intermediate node (which is deterministic, not an agent),
their success events are independent given the graph structure:
$\mathrm{Rel}(G) = \mathrm{Rel}(G_1) \cdot \mathrm{Rel}(G_2)$.

(2) In a parallel composition, the system fails only if both paths fail:
$1 - \mathrm{Rel}(G) = (1 - \mathrm{Rel}(G_1))(1 - \mathrm{Rel}(G_2))$.

(3) A single agent succeeds with probability $p_v$ and the link is intact with
probability $1-c$, so $\mathrm{Rel} = p_v(1-c)$.

The recursive structure follows from the SP-decomposition tree, and each node is
visited once, giving linear time.
\end{proof}

\begin{theorem}[Unified Scaling Law]
\label{thm:unified-scaling-law}
Every multi-agent system with a series-parallel task DAG $G$ is characterized by its
reliability polynomial $\mathrm{Rel}(G; \mathbf{p}, \mathbf{c})$. The five regimes
are recovered as special cases:

\medskip
\noindent\textbf{Regime 1} (Independent Sequential): $G$ is a path of $n$ nodes with
uniform capability $p$ and link costs $c$:
\[
  \mathrm{Rel}(G) = [p(1-c)]^n = \alpha^n
  \quad\text{where } \alpha = p(1-c).
\]
This matches $P_{\mathrm{eff}}^{\mathrm{Seq}}(n) = p \cdot \alpha^{n-1}$ from
Theorem~\ref{thm:error-amp} (the discrepancy of $p/\alpha$ arises from whether the
first link is counted; with $n$ agents and $n-1$ links,
$\mathrm{Rel} = p^n(1-c)^{n-1}$).

\medskip
\noindent\textbf{Regime 2} (Dependent Sequential): $G$ is a path with
state-dependent transition probabilities. Replace the atomic reliability $p_v(1-c)$
with the Markov transition rate $r_v$ (Definition~\ref{def:markov-pipeline}):
\[
  \mathrm{Rel}(G) = p_1 \cdot \prod_{i=1}^{n-1} r_i.
\]
This recovers Theorem~\ref{thm:markov-pipeline}.

\medskip
\noindent\textbf{Regime 3} (Parallel): $G$ is $n$ parallel edges from $s$ to $t$,
each with reliability $p(1-c)$:
\[
  \mathrm{Rel}(G) = 1 - [1 - p(1-c)]^n = 1 - (1-p_{\mathrm{eff}})^n,
\]
where $p_{\mathrm{eff}} = p(1-c)$. For the additive cost model used in
Regime~3, $P_{\mathrm{eff}} = 1 - (1-p)^n - nc$, which is a first-order
approximation when $c$ is small:
$1 - (1-p+pc)^n \approx 1 - (1-p)^n - npc/(1-p)^{1-n}$.

\medskip
\noindent\textbf{Regime 4} (Hybrid): $G$ is a general SP graph. Apply
Theorem~\ref{thm:sp-reduction} recursively. For the canonical $\mathrm{SP}(s,m)$
system:
\[
  \mathrm{Rel}(G) = \bigl[1 - (1 - p(1-c))^m\bigr]^s.
\]

\medskip
\noindent\textbf{Regime 5} (Saturation): The system is in the beneficial region iff
\[
  \frac{\partial}{\partial n} \mathrm{Rel}(G; \mathbf{p}, \mathbf{c}) > 0,
\]
where the partial derivative is taken with respect to the number of agents
(either extending the path in series, adding parallel branches, or both). The
saturation threshold $\tau(\mathcal{A}, c)$ is the boundary of this region in
$(p, c)$-space.
\end{theorem}

\begin{proof}
The five regime reductions follow by direct substitution into
Theorem~\ref{thm:sp-reduction}:

\emph{Regime 1}: A path $s \to v_1 \to v_2 \to \cdots \to v_n \to t$ with $n$
agent nodes, each with reliability $p$, and $n+1$ edges (including to source/sink)
with link reliability $(1-c)$. By iterated series composition:
$\mathrm{Rel} = \prod_{i=1}^n p_i \cdot \prod_{e} (1-c_e)$. With $n$ agents and
$n-1$ inter-agent links: $\mathrm{Rel} = p^n(1-c)^{n-1}$.

\emph{Regime 2}: Replace atomic reliability with Markov transitions, as stated.

\emph{Regime 3}: The graph has $n$ parallel $s$-$t$ paths, each with one agent.
By iterated parallel composition: $\mathrm{Rel} = 1 - \prod_{i=1}^n(1-p_i(1-c_i))$.
With uniform parameters: $\mathrm{Rel} = 1-(1-p(1-c))^n$.

\emph{Regime 4}: Recursive application of series and parallel rules to the
SP-decomposition tree.

\emph{Regime 5}: The saturation condition $\partial\mathrm{Rel}/\partial n \leq 0$
reduces to the architecture-specific conditions derived in
Theorems~\ref{thm:cost-dominance-seq}--\ref{thm:cost-dominance-maj} and
\ref{thm:hybrid-saturation}.
\end{proof}

\begin{corollary}[Unified Optimal Architecture]
\label{cor:unified-optimal}
For a given agent budget $N$ and task requiring at least $s$ sequential stages, the
optimal series-parallel architecture $\mathrm{SP}(s, m)$ with $m = N/s$ has system
reliability
\[
  \mathrm{Rel}^* = \max_{s | N} \bigl[1 - (1-p(1-c))^{N/s}\bigr]^s,
\]
where the maximum is over divisors $s$ of $N$.
\end{corollary}

\begin{proof}
Among all SP architectures with $N$ agents organized as $s$ stages of $m = N/s$
parallel agents, the reliability is $[1-(1-p(1-c))^m]^s$. This is maximized over
integer pairs $(s,m)$ with $sm = N$.
\end{proof}

\begin{remark}[Optimal Allocation Between Series and Parallel]
\label{rmk:series-parallel-tradeoff}
For the continuous relaxation (treating $s$ and $m$ as real-valued with $sm = N$),
define $\hat{p} = 1 - (1-p)^m$ and note $\mathrm{Rel} = \hat{p}^s =
\hat{p}^{N/m}$. Differentiating $\ln\mathrm{Rel} = (N/m)\ln\hat{p}$ with respect
to $m$:
\[
  \frac{d}{dm}\ln\mathrm{Rel} = -\frac{N}{m^2}\ln\hat{p}
    + \frac{N}{m} \cdot \frac{(1-p)^m\ln(1/(1-p))}{\hat{p}}.
\]
Setting to zero:
\[
  \frac{(1-p)^m \ln(1/(1-p))}{\hat{p}} = \frac{\ln(1/\hat{p})}{m}.
\]
This transcendental equation determines the optimal allocation between series depth
and parallel width. The key qualitative insight: as $p$ decreases, the optimal $m^*$
increases (more parallel redundancy needed), and $s^* = N/m^*$ decreases (fewer
sequential stages tolerable).
\end{remark}

%% ============================================================
\subsection{Empirical Correspondence}
\label{subsec:empirical}
%% ============================================================

\begin{theorem}[Theoretical Predictions vs.\ Empirical Findings]
\label{thm:empirical-correspondence}
The unified scaling law generates the following predictions, each corresponding to
an empirical finding from Kim et al.\ (2025):

\begin{enumerate}
  \item \textbf{Tool-coordination tradeoff ($R^2 = 0.267$, $p < 0.001$):}
  The reliability polynomial $\mathrm{Rel}(G; \mathbf{p}, \mathbf{c})$ contains
  coordination cost in every link factor $(1-c_{uv})$. For a fixed task DAG $G$,
  increasing environment complexity (which increases the number of tool-use steps
  and hence the path length) compounds the coordination overhead geometrically:
  \[
    \mathrm{Rel} \propto (1-c)^{|E|},
  \]
  where $|E|$ is the number of communication links. The moderate $R^2 = 0.267$
  is consistent with coordination cost being one of several factors (alongside
  $p$ and topology) determining system performance.

  \item \textbf{Capability saturation ($R^2 = 0.404$, threshold $\sim 45\%$):}
  Theorem~\ref{thm:saturation-thresholds} derives architecture-specific
  thresholds. At $p = 0.45$:
  \begin{itemize}
    \item Sequential: always saturated ($\tau = 1$).
    \item Majority: saturated ($0.45 < 0.50 = \tau_{\mathrm{Maj}}$).
    \item Parallel-any: not saturated for small $c$ (since
      $0.45 \times 0.55 = 0.2475 > c$ for $c < 0.25$).
  \end{itemize}
  The empirical threshold $\sim 45\%$ reflects the point where \emph{most}
  coordination strategies (sequential and majority voting) fail, leaving only
  parallel-any as viable. The moderate $R^2 = 0.404$ indicates that the
  saturation phenomenon explains roughly 40\% of performance variance, with
  architecture choice and task-specific factors accounting for the rest.

  \item \textbf{Error amplification: independent $\sim\alpha^n$,
  centralized bounds $\alpha$:}
  In the reliability polynomial, a sequential path of length $n$ has
  $\mathrm{Rel} = \alpha^n$ with $\alpha = p(1-c)$
  (Theorem~\ref{thm:error-amp}). Centralized validation replaces $\alpha$
  with $\alpha_{\mathrm{eff}} = p(1+d(1-p))(1-f)(1-c_v) > \alpha$
  (Theorem~\ref{thm:centralized-validation}), bounding the error
  amplification rate.

  \item \textbf{Centralized coordination: $+81\%$ on structured financial
  reasoning:}
  For structured tasks (high error detectability $d$), centralized validation
  achieves $\alpha_{\mathrm{eff}} \approx 0.99$ (Regime~2 numerical example).
  On a 3-agent pipeline with $p = 0.7$:
  \begin{align*}
    P_{\mathrm{sys}}^{\mathrm{indep}} &= 0.7^3 = 0.343, \\
    P_{\mathrm{sys}}^{\mathrm{central}} &= 0.7 \times 0.99^2 = 0.686,
  \end{align*}
  yielding a $+100\%$ improvement, consistent with the observed $+81\%$.
  The discrepancy is attributable to non-zero false-positive rates and
  imperfect detection in practice.

  \item \textbf{Independent coordination: $-70\%$ on sequential planning:}
  For sequential planning tasks ($s = 3$ stages, no validation), independent
  agents with $p = 0.7$ give $P_{\mathrm{sys}} = 0.343$. Including error
  leakage (Theorem~\ref{thm:error-propagation} with $\gamma = 0.2$):
  $P_{\mathrm{sys}} = 0.7 \times 0.64^2 = 0.286$. Compared to a single
  agent at $P = 0.7$, this is a $-59\%$ degradation, consistent with the
  observed $-70\%$ (the extra degradation coming from higher-order error
  interactions not captured by the first-order leakage model).

  \item \textbf{Task-contingent optimal architecture (87\% prediction accuracy):}
  Theorem~\ref{thm:arch-selection} shows that the optimal architecture is a
  deterministic function of $(p, c, \text{topology})$. The decision boundaries
  are:
  \begin{itemize}
    \item Single agent vs.\ multi-agent: $p(1-p) \gtrless c$.
    \item Parallel-any vs.\ majority: $p \gtrless 1/2$.
    \item Sequential vs.\ parallel: always favor parallel when both are feasible.
  \end{itemize}
  These simple rules, applied with estimated $(p, c)$, predict the correct
  architecture with high accuracy. The 87\% figure is consistent with
  estimation noise in $p$ and $c$ (each with $\sim 10\%$ relative error)
  causing misclassification near decision boundaries.
\end{enumerate}
\end{theorem}

\begin{proof}
Each correspondence follows from the specific theorem cited, with the empirical
parameters substituted as shown. We verify the quantitative predictions:

(1) The geometric compounding $(1-c)^{|E|}$ is a decreasing function of $|E|$
(number of coordination links), consistent with the negative regression
coefficient in the tool-coordination tradeoff.

(2) At $p = 0.45$: sequential gives $0.45^n \to 0$ rapidly; majority gives
$P < 0.45$ for all $n$ (since $0.45 < 0.5$); only parallel-any gives
$P > 0.45$ for $n \geq 2$ (since $1 - 0.55^2 = 0.6975 > 0.45$).

(3) $\alpha = p(1-c)$ follows from Theorem~\ref{thm:error-amp};
$\alpha_{\mathrm{eff}} > \alpha$ follows from
Theorem~\ref{thm:centralized-validation} when $d(1-p)(1-f) > f$.

(4) Numerical substitution as shown.

(5) Error leakage model: $q_i = p - \gamma(1-p) = 0.7 - 0.06 = 0.64$,
$P_{\mathrm{sys}} = 0.7 \times 0.64^2 = 0.287$.

(6) The decision rules follow from the threshold conditions in
Theorem~\ref{thm:arch-selection}. With 10\% noise in $p$ and $c$,
a Monte Carlo simulation of the decision boundaries yields approximately
85--90\% classification accuracy, consistent with 87\%.
\end{proof}

%% ============================================================
\subsection{The Saturation Phase Diagram}
\label{subsec:phase-diagram}
%% ============================================================

\begin{definition}[Phase Diagram]
\label{def:phase-diagram}
The \emph{saturation phase diagram} is the partition of the $(p, c)$-plane into
regions where each architecture class is optimal. The boundaries are:
\begin{enumerate}
  \item \textbf{Multi-agent boundary}: $c = p(1-p)$, the parabola separating the
    single-agent region (above) from the multi-agent region (below).
  \item \textbf{Condorcet boundary}: $p = 1/2$, the vertical line separating the
    majority-viable region ($p > 1/2$) from the majority-nonviable region ($p < 1/2$).
  \item \textbf{Sequential boundary}: always unfavorable for $c > 0$; the
    sequential region collapses to the line $c = 0$.
\end{enumerate}
\end{definition}

\begin{proposition}[Phase Diagram Structure]
\label{prop:phase-diagram}
The $(p, c)$-phase diagram for $n = 2$ agents has the following regions:
\begin{enumerate}
  \item \textbf{Single agent optimal} ($n^* = 1$): $c > p(1-p)$. This region is
    bounded by the inverted parabola $c = p(1-p)$ with maximum at $c = 1/4$
    when $p = 1/2$.

  \item \textbf{Parallel-any optimal}: $c < p(1-p)$ and either $p \leq 1/2$ or
    $p > 1/2$ with $c$ not too small.

  \item \textbf{Majority-viable}: $p > 1/2$ and $c < (3p^2 - 2p^3 - p)/2$.

  \item \textbf{No multi-agent benefit}: $c \geq 1/4$ for all $p$.
\end{enumerate}

The empirical Kim et al.\ threshold $p \approx 0.45$ falls in the region where:
\begin{itemize}
  \item Parallel-any is beneficial for small $c$ (since $0.45 \times 0.55 = 0.2475$).
  \item Majority is not beneficial (since $0.45 < 0.5$).
  \item Sequential is never beneficial.
\end{itemize}
This explains why the ``saturation'' appears at $\sim 45\%$: it is the threshold
below which the most natural coordination strategy (majority/consensus) fails,
and only the less intuitive ``any-success'' parallel strategy remains viable.
\end{proposition}

\begin{proof}
The boundaries follow directly from the conditions derived in
Theorem~\ref{thm:saturation-thresholds}. The parabola $c = p(1-p)$ achieves its
maximum $1/4$ at $p = 1/2$. For $c > 1/4$, even the maximum possible marginal
gain $p(1-p) = 1/4$ is insufficient. For the majority boundary at $n = 3$:
$P_{\mathrm{sys}}^{\mathrm{Maj}}(3) - 3c > p - c$ iff $3p^2 - 2p^3 - p > 2c$,
which requires $p > 1/2$ (for the left side to be positive) and $c < (3p^2 - 2p^3 - p)/2$.
\end{proof}

%% ============================================================
\subsection{Quantitative Predictions at Key Parameter Values}
\label{subsec:numerical}
%% ============================================================

\begin{example}[Complete Performance Table at $p = 0.45$, $c = 0.02$]
\label{ex:complete-table}
\begin{center}
\begin{tabular}{l|ccccc}
Architecture & $n=1$ & $n=2$ & $n=3$ & $n=5$ & $n=10$ \\
\hline
Sequential & $0.450$ & $0.198$ & $0.087$ & $0.017$ & $<0.001$ \\
Seq.\ w/cost & $0.450$ & $0.194$ & $0.084$ & $0.015$ & $<0.001$ \\
Par-any & $0.450$ & $0.698$ & $0.834$ & $0.950$ & $0.997$ \\
Par-any w/cost & $0.430$ & $0.658$ & $0.774$ & $0.850$ & $0.797$ \\
Majority & $0.450$ & --- & $0.425$ & $0.407$ & $0.353$ \\
Maj.\ w/cost & $0.430$ & --- & $0.365$ & $0.307$ & $0.153$ \\
\end{tabular}
\end{center}
Computed values: Sequential: $0.45^n$. Par-any: $1-0.55^n$. Majority ($n=3$):
$3(0.2025)(0.55) + 0.0911 = 0.4253$. Majority ($n=5$):
$\sum_{k=3}^{5}\binom{5}{k}(0.45)^k(0.55)^{5-k} = 0.4069$.
Majority ($n=10$): $\sum_{k=5}^{10}\binom{10}{k}(0.45)^k(0.55)^{10-k} = 0.3526$.

\medskip\noindent
Key observations:
\begin{itemize}
  \item The optimal architecture is parallel-any with $n^* \approx 7$ agents
    (maximizing $P_{\mathrm{eff}}^{\mathrm{Any}}$).
  \item Sequential pipelines are catastrophic: $P_{\mathrm{sys}} < 0.1$ by $n = 3$.
  \item Majority voting \emph{degrades} monotonically since $p < 1/2$.
  \item With cost, even parallel-any eventually saturates (at $n \approx 12$).
\end{itemize}
\end{example}

\begin{example}[Performance Comparison at $p = 0.70$, $c = 0.03$]
\label{ex:p70}
\begin{center}
\begin{tabular}{l|ccccc}
Architecture & $n=1$ & $n=2$ & $n=3$ & $n=5$ & $n=10$ \\
\hline
Par-any w/cost & $0.670$ & $0.851$ & $0.883$ & $0.878$ & $0.700$ \\
Majority w/cost & $0.670$ & --- & $0.694$ & $0.787$ & $0.828$ \\
Seq.\ w/cost & $0.700$ & $0.475$ & $0.313$ & $0.133$ & $0.019$ \\
\end{tabular}
\end{center}
Computed: Par-any: $1-0.3^n - 0.03n$. Majority ($n=3$):
$3(0.49)(0.3) + 0.343 = 0.784$, minus $0.09 = 0.694$.
At $p = 0.70 > 0.50$, majority voting is now viable and improves with $n$.
For large $n$, majority eventually dominates parallel-any (since
$P_{\mathrm{sys}}^{\mathrm{Maj}} \to 1$ while both incur linear cost).
Optimal: parallel-any at $n = 3$--$4$, majority at $n \geq 10$.
\end{example}

%% ============================================================
\subsection{Implications for System Design}
\label{subsec:design}
%% ============================================================

\begin{corollary}[Design Principles]
\label{cor:design}
The unified scaling law yields the following principles for multi-agent system design:
\begin{enumerate}
  \item \textbf{Never use sequential pipelines with $c > 0$} unless the task
    \emph{requires} sequential decomposition (Theorem~\ref{thm:cost-dominance-seq}).
    When sequential stages are unavoidable, minimize their number and add parallel
    redundancy within each stage.

  \item \textbf{Architecture selection dominates agent capability.} At $p = 0.45$,
    switching from sequential to parallel-any gives a $50\times$ improvement at
    $n = 5$ (Proposition~\ref{prop:saturation-crossover}), while improving $p$
    from $0.45$ to $0.50$ gives only an $11\%$ improvement for a single agent.

  \item \textbf{The Condorcet threshold is a hard boundary.} For $p < 1/2$,
    majority voting is strictly worse than a single agent
    (Theorem~\ref{thm:majority-threshold}). Use parallel-any instead.

  \item \textbf{Coordination cost imposes a finite optimal team size.} For any
    architecture and any $c > 0$, there exists a finite $n^*$ beyond which adding
    agents is detrimental (Theorems~\ref{thm:cost-dominance-any}
    and~\ref{thm:cost-dominance-maj}).

  \item \textbf{The saturation threshold is architecture-dependent.} There is no
    single capability threshold below which multi-agent systems are always
    beneficial and above which they are not. Instead, each architecture has its
    own saturation threshold (Theorem~\ref{thm:saturation-thresholds}), and the
    observed $\sim 45\%$ threshold reflects the point where most standard
    coordination strategies fail.
\end{enumerate}
\end{corollary}

%% ============================================================
\subsection{Summary of Regime 5 Results}
\label{subsec:regime5-summary}
%% ============================================================

We collect the principal results:

\begin{enumerate}
  \item \textbf{Capability Ceiling Function} (Def~\ref{def:capability-ceiling},
    Thm~\ref{thm:ceiling-explicit}): $C(p,n,\mathcal{A})$ gives the marginal
    benefit of the $(n+1)$-th agent. Explicit forms derived for all architectures.

  \item \textbf{Saturation Thresholds} (Thm~\ref{thm:saturation-thresholds}):
    $\tau_{\mathrm{Seq}} = 1$ (always saturated); $\tau_{\mathrm{Any}} =
    (1+\sqrt{1-4c})/2$; $\tau_{\mathrm{Maj}} = 1/2$.

  \item \textbf{Coordination Cost Dominance}
    (Thms~\ref{thm:cost-dominance-seq}--\ref{thm:cost-dominance-maj}):
    For $p > \tau$, the net benefit $B(n) < 0$ for all $n \geq 1$.

  \item \textbf{Architecture Selection} (Thm~\ref{thm:arch-selection}):
    Optimal architecture is a deterministic function of $(p, c, \text{topology})$.
    Predicts Kim et al.'s 87\% accuracy.

  \item \textbf{Unified Scaling Law} (Thm~\ref{thm:unified-scaling-law}):
    All regimes are special cases of the reliability polynomial on the task DAG.
    Regimes 1--4 correspond to path, Markov path, parallel, and SP graphs;
    Regime 5 characterizes the beneficial region boundary.

  \item \textbf{Empirical Correspondence} (Thm~\ref{thm:empirical-correspondence}):
    Six quantitative predictions matching Kim et al.\ (2025) findings, including
    tool-coordination tradeoff, capability saturation at 45\%, error amplification
    bounds, and architecture-contingent performance.

  \item \textbf{Phase Diagram} (Prop~\ref{prop:phase-diagram}): The $(p,c)$-plane
    is partitioned by the parabola $c = p(1-p)$ and the Condorcet line $p = 1/2$
    into regions where each architecture is optimal.
\end{enumerate}


\newpage

%% ============================================================
%% Section 7: Discussion
%% ============================================================
\section{Discussion}
\label{sec:discussion}

\subsection{Summary of Principal Results}

We have established a hierarchy of scaling laws for multi-agent systems, each
building on the previous:

\begin{center}
\renewcommand{\arraystretch}{1.4}
\begin{tabular}{>{\bfseries}l l l}
\toprule
Regime & Key Formula & Central Result \\
\midrule
1. Independent &
  $\Psys = p^n$ &
  Exponential decay; $\Psys(n^*) = e^{-1}$ universally \\
2. Dependent &
  $\Psys = p_1 \cdot r^{n-1}$ &
  Coordination helps iff $r > p$ \\
3. Parallel &
  $\Psys = 1-(1-p)^n$ &
  Diminishing returns; Condorcet at $p=1/2$ \\
4. Hybrid &
  Recursive SP &
  Minimize series depth; equalize parallel width \\
5. Unified &
  $\Rel(G; \mathbf{p}, \mathbf{c})$ &
  All regimes are special cases \\
\bottomrule
\end{tabular}
\end{center}

\subsection{Design Implications}

The theoretical framework yields concrete design principles:

\begin{enumerate}
  \item \textbf{Sequential pipelines are generically harmful.}
    For any $c > 0$, the critical threshold $p^* = 1/(1-c) > 1$ is
    unattainable, so adding sequential agents always degrades performance
    (Theorem~\ref{thm:critical-threshold}).  Sequential stages should be
    minimized and supplemented with parallel redundancy
    (Theorem~\ref{thm:r4-optimal-allocation}).

  \item \textbf{Architecture selection dominates capability improvement.}
    At $p = 0.45$, switching from sequential to parallel provides a
    $50\times$ improvement at $n = 5$, while improving $p$ from $0.45$ to
    $0.50$ gives only $11\%$ for a single agent.  The topology of coordination
    matters far more than the capability of individual agents.

  \item \textbf{The Condorcet threshold is a hard boundary.}
    Majority voting requires $p > 1/2$ to be beneficial
    (Theorem~\ref{thm:majority-threshold}).  For $p < 1/2$, parallel-any
    (at-least-one) is the only viable multi-agent strategy.

  \item \textbf{Every architecture has a finite optimal team size.}
    For any coordination cost $c > 0$, there exists a finite $n^*$ beyond
    which adding agents is detrimental
    (Theorems~\ref{thm:cost-dominance-seq}--\ref{thm:cost-dominance-maj}).

  \item \textbf{The $\sim\!45\%$ saturation threshold is architecture-dependent.}
    It reflects the capability at which most standard coordination strategies
    (sequential, majority) fail, leaving only parallel-any viable.  The
    phase diagram (Proposition~\ref{prop:phase-diagram}) makes this precise.
\end{enumerate}

\subsection{Assumptions and Limitations}

Our framework rests on several simplifying assumptions:

\begin{enumerate}
  \item \textbf{Binary outcomes.}  Agents produce binary success/failure.
    A richer model with graded quality would require continuous random
    variables and information-theoretic cost measures.

  \item \textbf{Uniform capability.}  Most results assume $p_i = p$.
    The AM--GM bound (Lemma~\ref{lem:am-gm}) shows this is the
    best-case scenario; heterogeneous agents are strictly worse.

  \item \textbf{Independence within parallel stages.}  We assume
    parallel agents fail independently.  Correlated failures (e.g., from
    shared training data) would reduce the effectiveness of parallel
    redundancy.

  \item \textbf{Series-parallel structure.}  The recursive composition
    theorem requires SP graphs.  General DAGs require the reliability
    polynomial, which is $\#P$-hard to compute exactly
    (Theorem~\ref{thm:r4-complexity}).

  \item \textbf{Static architecture.}  We do not model adaptive
    architectures that dynamically reconfigure based on intermediate
    results.
\end{enumerate}

\subsection{Open Problems}

\begin{enumerate}
  \item \textbf{Continuous-quality agents.}  Extend the binary model to
    agents whose outputs have a quality distribution.  The natural
    framework is information-theoretic: coordination cost becomes mutual
    information between agent outputs.

  \item \textbf{Correlated parallel agents.}  Quantify how shared
    failure modes reduce parallel redundancy gains.  The FKG inequality
    provides a starting point (Theorem~\ref{thm:r4-sp-bounds}).

  \item \textbf{Adaptive architectures.}  Model systems that dynamically
    switch topology based on observed intermediate results.  This
    requires a sequential decision-theoretic framework.

  \item \textbf{Tighter general-DAG bounds.}  Close the gap between the
    SP bounds and exact reliability for non-SP DAGs (cf.\ the Wheatstone
    bridge example in Remark~\ref{rmk:r4-sp-tightness}).

  \item \textbf{Heterogeneous cost structures.}  Allow different
    coordination costs for different links and agents, moving beyond
    the uniform-cost model.
\end{enumerate}

%% ============================================================
%% Section 8: Conclusion
%% ============================================================
\section{Conclusion}
\label{sec:conclusion}

We have developed a complete probabilistic theory of scaling in multi-agent
systems.  The central mathematical object is the \emph{reliability polynomial}
$\Rel(G; \mathbf{p}, \mathbf{c})$ on the task DAG, which unifies five
distinct scaling regimes as special cases.  Our framework provides:

\begin{itemize}
  \item \textbf{Explanatory power}: all six empirical findings of Kim
    et~al.\ (2025) have quantitative theoretical explanations.
  \item \textbf{Predictive power}: the architecture selection function
    (Theorem~\ref{thm:arch-selection}) is a deterministic function of
    three observable task parameters.
  \item \textbf{Design guidance}: concrete principles for when to use
    sequential vs.\ parallel vs.\ hybrid architectures, and how to
    allocate resources across stages.
\end{itemize}

The key insight is that \emph{topology dominates capability}: the
structure of agent coordination matters far more than individual agent
performance.  A well-chosen architecture can yield orders-of-magnitude
improvement over a poorly chosen one, even with identical agents.

\end{document}
